\PassOptionsToPackage{svgnames}{xcolor}
\documentclass{article}
\usepackage{fontspec,symbats3,array,parskip,fancyhdr,setspace,relsize}
\usepackage[backend=biber,style=authoryear]{biblatex}
\usepackage{enumitem,xcolor,multicol,menukeys,sectsty,metalogo}
\usepackage[normalem]{ulem}
\usepackage{fancybox}
%\usepackage{draftwatermark}
%\SetWatermarkColor{LightSteelBlue}
\addbibresource{symbats3.bib}
\ULdepth=0pt
\makeatletter
\UL@height=.4pt
\renewcommand\uline{%
  \bgroup\markoverwith{\textcolor{Red}{\rule[\ULdepth]{2pt}{\UL@height}}}\ULon}
\allsectionsfont{\sffamily}
\makeatother
\setmainfont{Alegreya}
\setsansfont{Alegreya Sans}
\setmonofont[Scale=MatchLowercase]{InconsolataN}
\usepackage[a4paper,margin=1cm,top=2cm,bottom=2cm]{geometry}
\def\,{,}
\def\<{\texttt{<}}
\def\|{\textbar}
\def\>{\texttt{>}}
\setcounter{tocdepth}{1}
\newcommand{\vstrut}{\vrule height1.2em width0pt}
\pagestyle{fancy}
\rhead{\scriptsize\leftmark}
\lhead{\scriptsize\rightmark}
\renewcommand{\abstractname}{Summary}
\usepackage{hyperref}
\hypersetup{colorlinks=true,linkcolor=Crimson,urlcolor=MediumBlue,citecolor=Green}
\AtBeginDocument{\raggedright\thispagestyle{empty}}
\begin{document}
\title{Symbats 3.0 font macros for \LaTeX}
\author{Peter Flynn\\Silmaril Consultants\\\texttt{peter@silmaril.ie}}
\maketitle
\begin{center}
\begin{minipage}{.75\textwidth}
\begin{abstract}\onehalfspacing\raggedright\parindent1em\noindent
{S\smaller YMBATS} is a Neopagan dingbats typeface family that includes
various religious, astrological and other symbols, plus support for
Runic and Ogham script. The designer, Feòrag, has kindly released the
fonts under the SIL Open Font License Version 1.1 (26 February 2007).

This \LaTeX\ package (\textsf{symbats}) is released under the
\LaTeX\ Project Public License (LPPL), version 1.3c and is available
from the Comprehensive \TeX\ Archive Network (CTAN).

\subsubsection*{Current status}
\begin{enumerate}
  \item Names for use as commands in \LaTeX\ have been specified for
    all identified glyphs
  \item Names for glyphs positioned for older versions of the fonts
    have been prefixed `old' where this was identifiable;
  \item Names for Runic are prefixed `RU' and names for Ogham are
    prefixed `OG';
  \item Names which otherwise appear to be duplicates have been
    prefixed `dup'.
\end{enumerate}
\subsubsection*{Errors and omissions}
The maintainer of this package is Peter Flynn
\href{mailto:peter@silmaril.ie}{\nolinkurl{peter@silmaril.ie}}. Please
notify any errors or omissions, especially corrections to the naming,
to the maintainer.
\end{abstract}
\tableofcontents
\subsection*{History}
\begin{itemize}[nosep]
  \item[0.9] 2022-06-21 Edits and corrections to descriptions
  \item[0.8] 2022-01-18 Baseline adjustment option added,
    documentation updated
  \item[0.7] 2022-01-11 Last duplicated names fixed, .sty file created
    manually
  \item[0.6] 2022-01-08 Names edited and normalised, out for checking
  \item[0.5] 2022-01-06 Got feedback, edited names as suggested,
    created Makefile
  \item[0.4] 2022-01-04 Sent name list to Feòrag for checking
  \item[0.3] 2022-01-03 Created CSV and \LaTeX\ files
  \item[0.2] 2022-01-02 First pass at constructing names for the
    glyphs
  \item[0.1] 2021-12-20 Downloaded font, extracted list of glyph
    codepoints
\end{itemize}
\end{minipage}
\end{center}
\clearpage
\begin{multicols*}{2}
\large
\section{Symbats 3.0}
Symbats is a Neopagan dingbats font which has been developed by Feòrag
over a couple of decades, and recently re-released in an updated
version announced in a Twitter
post\footnote{\url{https://twitter.com/Feorag/status/1473048711601266690?s=20}}
on Winter Solstice Eve 2021.

The font comes in three weights, and has many additions and
redesigns. The files are distributed in OpenType format
(\texttt{.otf}) from the author's
website.\footnote{\url{https://www.feorag.com/freestuff/symbats.html}}\label{intro}

\subsection{\LaTeX\ version}

For the \LaTeX\ version, I have assigned control sequences
(\LaTeX\ commands) to all the glyphs. They are listed in this document
three times, in different orders for ease of reference:
\begin{enumerate}[nosep]
\item index order of Unicode codepoint;
\item alphabetical order of \LaTeX\ name;
\item alphabetical order of Unicode name.
\end{enumerate}

\subsection{Numeric character positions}\label{numpos}

For historical reasons, some of the glyphs are in the code positions
of ASCII printable characters (U+0021 to U+007E), where they were
originally placed for ease of typing in pre-Unicode days.

For compatibility with older versions, these glyphs are still in
place, but in cases where there are now real Unicode codepoints for the
exact characters by name, the glyph is duplicated at the new position
and the glyph at the original location has its name prefixed with
`\texttt{old}'. To the best of my understanding, the two glyphs are
identical in each case.

These duplications are noted in the following tables in the
\textbf{Dupes} column with a clickable link to the other position
(note that all these links go to the first listing, the index order of
codepoint, not to the listing in which you clicked).

For example, Earth (\oldearth) was originally at U+003C, but
Unicode now defines an actual codepoint named `Earth' at U+2641, so
the glyph is named in that location as \verb+\earth+ and the
original at U+003C is named \verb+\oldearth+; but they are identical.

\subsection{Naming}

In choosing the names, the methodology used was:

\begin{enumerate}

\item If the glyph has an established name (Sun, Earth, Pentagram)
  that name is used in all lowercase, eg \verb+\earth+;

\item Glyphs with multiple variants (moon, new moon, waxing moon,
  waning moon, etc) are named with the principal component first, eg
  \verb+\moonnew+, \verb+\moonfull+, \verb+\moonwaningoutline+,
  \verb+\moonwaningsolid+, etc.

  These can best be compared in the Unicode name order listing in
  \hyperref[unicodename]{\S\ref{unicodename}} starting on
  p.\pageref{unicodename}.

\item Variations in \textbf{pentagrams} include:
  \begin{itemize}[nosep]
  \item \textbf{rough} (sketched lines);
  \item \textbf{inverted} (abbreviated to \texttt{inv} to save typing);
  \item \textbf{interlaced} (lines pass in front and behind others);
  \item handedness of the interlacing (\textbf{Right} \textit{vs}
    \textbf{Left}) using capital \texttt{R} and \texttt{L};
  \item \textbf{circled} (ring around) \textit{vs} \textbf{solid}
    (reversed out).
  \end{itemize}
  For the first three (rough, inverted, and interlaced) there is no
  explicit opposite in the name, so the complement of
  \verb+\pentagramroughcircled+ is just \verb+\pentagramcircled+:
  there is no `smooth' in the name.

  Similarly, upright symbols are not
  named as such, only the inverted versions get an `\texttt{inv}' in
  their name.

  Where there is an obvious complement, this is linked in the
  \textbf{Compl} column in a similar manner to the \textbf{Dupes}
  method described in \hyperref[numpos]{\S\ref{numpos}}.

\item Variations in \textbf{sun} and \textbf{moon} include:
  \begin{itemize}[nosep]
  \item glyph in \textbf{outline} or \textbf{solid} (ie empty or filled);
  \item sun — number of \textbf{spokes} [is this the correct term? rays?];
  \item sun — \textbf{corona} (external ring);
  \item sun — \textbf{hub} (centre ring);
  \item moon — phase (new, first quarter, third quarter, full, waxing,
    waning).
  \end{itemize}
  As with the pentagrams, the absence of an effect is signalled by
  the absence of the term in the name.

\item Ogham letter names are prefixed with `\texttt{OG}';

\item Runic letter names are prefixed with `\texttt{RU}';

\item For any existing glyph for which a formal Unicode codepoint now
  exists, the name in the original location is prefixed by
  `\texttt{old}' as described in \hyperref[numpos]{\S\ref{numpos}};

\item A few names are not known, or not readily identifiable, or lost
  to the mists of time. These have been guessed at, and I would
  welcome suggestions.
  
\end{enumerate}

\section{Using the \textit{symbats3} package}

Download the fonts (see \S\ref{intro}) and install them into
  your computer according to your operating system's instructions
  (see instructions at the end of this page).

\subsection{Package installation}

\subsubsection{Normal (automated) installation}

If you are using (eg) MiK\TeX\ or another system with automated font
installation, just invoking the package in a document is enough to
make it install automatically from CTAN, eg:

\verb+\usepackage{symbats3}+

\subsubsection{Manual installation}

If you have to (or want to) install the package manually:

\begin{enumerate}
  
\item Download the package zip file from CTAN.

\item Unzip the file into your Personal \TeX\ Folder (see
  page\thinspace\pageref{ptf}) 

\item Create the subdirectories \verb+tex/latex/symbats3/+ and
  \verb+doc/latex/symbats3/+ within
    your Personal \TeX\ Folder \verb+texmf+.
    
\item Move the two package files \verb+symbats3.sty+ and
  \verb+symbats3.def+ into the new \LaTeX\ subdirectory
  \verb+texmf/tex/latex/symbats3/+.

\item Move the remaining package files into the new documentation
    subdirectory \verb+texmf/doc/latex/symbats3/+.

\item Add the line \verb+\usepackage{symbats3}+ to the Preamble of
  your \LaTeX\ document.

  \end{enumerate}

Note that the \textsf{supertramp} package \textsf{\smaller REQUIRES}
the \textsf{fontspec} package for Unicode and OpenType fonts
\parencite{fontspec}. The \textsf{supertramp} package \textsf{\smaller
  CANNOT} be used with the old \textsf{fontenc} and \textsf{inputenc}
packages, which means you \textsf{\smaller MUST} use \XeLaTeX\ or
  \LuaLaTeX\ as your processor, not \textit{pdflatex} or plain
  \textit{latex}.

\subsection{Options}

There is only one option to this version of the package:
\verb+[descenders]+, which lowers the height of the glyphs to align with
the descenders of the surrounding text.

By default, the glyphs print with their baseline aligned with the
baseline of the surrounding text, eg

\uline{\hbox{\sffamily\huge Baseline: \pentagram}}

In many cases the glyph fits better if it is aligned instead with the
bottom of the descenders (q, y, p, j, g) of the surrounding text, eg

\uline{\hbox{\sffamily\huge Descenders (qypjg): \pentagram[-.5ex]}}

The package option \verb+[descenders]+ will automatically calculate the
depth of the descenders for the current font and lower the glyph to
align with them, eg

\verb+\usepackage[descenders]{symbats3}+

In addition, there is an optional argument to every glyph command to
specify a different height or depth, eg

\verb+\pentagraminvsolid[-.666ex]+

will produce

\uline{\hbox{\sffamily\huge Manual: \pentagraminvsolid[-.666ex]}}

\bigskip\hrule\bigskip

\section*{Note on required features}
In this document, the keywords
    {\sffamily {\smaller MUST}}, {\sffamily {\smaller MUST NOT}}, {\sffamily {\smaller REQUIRED}},
    {\sffamily {\smaller SHALL}}, {\sffamily {\smaller SHALL NOT}}, {\sffamily {\smaller SHOULD}},
    {\sffamily {\smaller SHOULD NOT}},
    {\sffamily {\smaller RECOMMENDED}},
    {\sffamily {\smaller MAY}}, and
    {\sffamily {\smaller OPTIONAL}} have a specific
    meaning when shown in {\sffamily {\smaller THIS TYPESTYLE}}, and
    {\sffamily {\smaller MUST}} be interpreted as described in
    RFC 2119 \parencite{rfc2119}.
When shown in another typestyle, these words keep their conventional
    contextual degree of meaning.\par

\printbibliography


\section*{Installing OpenType™ and TrueType™ fonts}

  \begin{description}
    \item[Linux:] Run the \textsf{Font Manager} app, click
      \keys{\texttt{+}} and select the font files to install.
    \item[Mac OS X] Run the \textsf{Font Book} app, click the
      \textsf{Add} button in the Font Book toolbar, locate and select
      a font, then click \textsf{Open}. Drag the font file to the Font
      Book app icon in the Dock. Double-click the font file in the
      Finder, then click \textsf{Install Font} in the dialog that
      appears.
    \item[Windows:] Run \textsf{Directory Explorer} (what used to be
      called \textsf{My Computer}), right-click on
      each font file where you downloaded it, and select \textsf{Install}.
  \end{description}


  \begingroup
  \fboxsep1em
  \begin{Sbox}
    \begin{minipage}{.9\hsize}
      \raggedright\sffamily
      Your Personal \TeX\ Folder should normally be something like this;\\
      create it now if it does not already exist:
      \par\smallskip
      \begin{description}[noitemsep,leftmargin=1em,font=\sffamily]
      \item[Linux:] \verb+~/texmf+
      \item[Mac OS\thinspace X:]
        \verb+~/Library/texmf+
      \item[Windows:] (\textsf{Win95–XP})
        \verb+C:\texmf\tex\latex\symbats3+\\
        (\textsf{Win 7–11})
        \verb+Computer\System\Users\+\texttt{\itshape yourname}\verb+\texmf+
        \\[2mm]
        \textbf{Note} Windows MiK\TeX\ users (only) \textsf{\smaller
          MUST} add their Personal \TeX\ Folder to their MiK\TeX\ root
        when it is first created; and \textsf{\smaller MUST} refresh their
        Font Name DataBase (FNDB) utility after installation: see the
        link below. No such action is needed for other systems.
      \end{description}
      \par\smallskip
      See the online documentation at
      \url{http://latex.silmaril.ie/formattinginformation/personal.html}
      for more information about your
      Personal \TeX\ Folder and about how to use it in MiK\TeX\ \parencite{fi}.
    \end{minipage}
  \end{Sbox}
  \fbox{\TheSbox}
  \endgroup\label{ptf}

\end{multicols*}

  \clearpage
\section{Symbats3 in index order of Unicode codepoint}\label{codepoint}
\begin{tabular}{@{}
  >{\ttfamily}l@{}
  >{\sffamily}c@{\hspace{2mm}}
  >{\sffamily\raggedright\vstrut}p{.3\textwidth}
  >{\Large\symbats}c
  >{\ttfamily\raggedright\arraybackslash\textbackslash}m{.29\textwidth}
  >{\ttfamily}l@{\hspace{2mm}}
  >{\ttfamily}l@{}}
\normalfont\bfseries Code&
\normalfont\bfseries Char&
\normalfont\bfseries Unicode name&
\normalfont\normalsize\bfseries Glyph&
\multicolumn1{>{\normalfont\bfseries}l}{\LaTeX\ name}&
\normalfont\bfseries Dupes&
\normalfont\bfseries Compl\\\hline\vrule height1.75em width0pt
\hypertarget{U+0021}{U+0021}
&!&Exclamation Mark&\moonwaxfullwaneoutline&moonwaxfullwaneoutline&
&
\\
\hypertarget{U+0022}{U+0022}
&“&Quotation Mark&\valknutsolid&valknutsolid&
&
\\
\hypertarget{U+0023}{U+0023}
&\#&Number Sign&\oldgemini&oldgemini&
\hyperlink{U+264A}{U+264A}
&
\\
\hypertarget{U+0024}{U+0024}
&\$&Dollar Sign&\moonfirstquarteroutline&moonfirstquarteroutline&
&
\\
\hypertarget{U+0025}{U+0025}
&\%&Percent Sign&\moonnew&moonnew&
&
\\
\hypertarget{U+0026}{U+0026}
&\&&Ampersand&\oldmoonwaningoutline&oldmoonwaningoutline&
\hyperlink{U+263E}{U+263E}
&
\\
\hypertarget{U+0027}{U+0027}
&'&Apostrophe&\lightning&lightning&
&
\\
\hypertarget{U+0028}{U+0028}
&(&Left Parenthesis&\moonfirstquartersolid&moonfirstquartersolid&
&
\\
\hypertarget{U+0029}{U+0029}
&)&Right Parenthesis&\moonfull&moonfull&
&
\\
\hypertarget{U+002A}{U+002A}
&*&Asterisk&\moonwaxingsolid&moonwaxingsolid&
&
\\
\hypertarget{U+002B}{U+002B}
&+&Plus Sign&\moonwaningsolid&moonwaningsolid&
&
\\
\hypertarget{U+002C}{U+002C}
&\,&Comma&\oldceres&oldceres&
\hyperlink{U+26B3}{U+26B3}
&
\\
\hypertarget{U+002D}{U+002D}
&-&Hyphen-Minus&\sunthreedoublespokes&sunthreedoublespokes&
&
\\
\hypertarget{U+002E}{U+002E}
&.&Full Stop&\sulphur&sulphur&
&
\\
\hypertarget{U+002F}{U+002F}
&/&Solidus&\crossedstaff&crossedstaff&
&
\\
\hypertarget{U+0030}{U+0030}
&0&Digit Zero&\sunfourspokes&sunfourspokes&
&
\\
\hypertarget{U+0031}{U+0031}
&1&Digit One&\sunwithrays&sunwithrays&
&
\\
\hypertarget{U+0032}{U+0032}
&2&Digit Two&\sunwithraysfourspokes&sunwithraysfourspokes&
&
\\
\hypertarget{U+0033}{U+0033}
&3&Digit Three&\oldsuncoronafourspokes&oldsuncoronafourspokes&
\hyperlink{U+0038}{U+0038}
&
\\
\hypertarget{U+0034}{U+0034}
&4&Digit Four&\octile&octile&
&
\\
\hypertarget{U+0035}{U+0035}
&5&Digit Five&\sun&sun&
&
\\
\hypertarget{U+0036}{U+0036}
&6&Digit Six&\sunfourdoublespokes&sunfourdoublespokes&
&
\\
\hypertarget{U+0037}{U+0037}
&7&Digit Seven&\suneightspokes&suneightspokes&
&
\\
\hypertarget{U+0038}{U+0038}
&8&Digit Eight&\suncoronafourspokes&suncoronafourspokes&
\hyperlink{U+0033}{U+0033}
&
\\
\hypertarget{U+0039}{U+0039}
&9&Digit Nine&\suncoronahubfourspokes&suncoronahubfourspokes&
&
\\
\hypertarget{U+003A}{U+003A}
&:&Colon&\valknutoutline&valknutoutline&
&
\\
\hypertarget{U+003B}{U+003B}
&;&Semicolon&\stareightarrows&stareightarrows&
&
\\
\hypertarget{U+003C}{U+003C}
&\<&Less-Than Sign&\oldearth&oldearth&
\hyperlink{U+2641}{U+2641}
&
\\
\hypertarget{U+003D}{U+003D}
&=&Equals Sign&\suncoronasixspokes&suncoronasixspokes&
&
\\
\hypertarget{U+003E}{U+003E}
&\>&Greater-Than Sign&\recipe&recipe&
&
\\
\hypertarget{U+003F}{U+003F}
&?&Question Mark&\celtictrefoilknotopen&celtictrefoilknotopen&
&
\\
\hypertarget{U+0040}{U+0040}
&@&Commercial At&\moonwaxfullwanesolid&moonwaxfullwanesolid&
&
\\
\hypertarget{U+0041}{U+0041}
&A&Latin Capital Letter A&\whiteknifesouth&whiteknifesouth&
&
\\
\hypertarget{U+0042}{U+0042}
&B&Latin Capital Letter B&\oldsextile&oldsextile&
\hyperlink{U+26B9}{U+26B9}
&
\\
\hypertarget{U+0043}{U+0043}
&C&Latin Capital Letter C&\oldconjunction&oldconjunction&
\hyperlink{U+260C}{U+260C}
&
\\
\hypertarget{U+0044}{U+0044}
&D&Latin Capital Letter D&\whiteknifeeast&whiteknifeeast&
&
\\
\end{tabular}
\par
\begin{tabular}{@{}
  >{\ttfamily}l@{}
  >{\sffamily}c@{\hspace{2mm}}
  >{\sffamily\raggedright\vstrut}p{.3\textwidth}
  >{\Large\symbats}c
  >{\ttfamily\raggedright\arraybackslash\textbackslash}m{.29\textwidth}
  >{\ttfamily}l@{\hspace{2mm}}
  >{\ttfamily}l@{}}
\normalfont\bfseries Code&
\normalfont\bfseries Char&
\normalfont\bfseries Unicode name&
\normalfont\normalsize\bfseries Glyph&
\multicolumn1{>{\normalfont\bfseries}l}{\LaTeX\ name}&
\normalfont\bfseries Dupes&
\normalfont\bfseries Compl\\\hline\vrule height1.75em width0pt
\hypertarget{U+0045}{U+0045}
&E&Latin Capital Letter E&\pentagraminvintLcircled&pentagraminvintLcircled&
&
\hyperlink{U+00C9}{U+00C9}
\\
\hypertarget{U+0046}{U+0046}
&F&Latin Capital Letter F&\whiteknifewest&whiteknifewest&
&
\\
\hypertarget{U+0047}{U+0047}
&G&Latin Capital Letter G&\mazesquarebase&mazesquarebase&
&
\\
\hypertarget{U+0048}{U+0048}
&H&Latin Capital Letter H&\mazeround&mazeround&
&
\\
\hypertarget{U+0049}{U+0049}
&I&Latin Capital Letter I&\pentagraminvintRsolid&pentagraminvintRsolid&
&
\hyperlink{none}{none}
\\
\hypertarget{U+004A}{U+004A}
&J&Latin Capital Letter J&\rockartgoddess&rockartgoddess&
&
\\
\hypertarget{U+004B}{U+004B}
&K&Latin Capital Letter K&\rockarthorse&rockarthorse&
&
\\
\hypertarget{U+004C}{U+004C}
&L&Latin Capital Letter L&\labrys&labrys&
&
\\
\hypertarget{U+004D}{U+004D}
&M&Latin Capital Letter M&\dupsunfourspokes&dupsunfourspokes&
&
\\
\hypertarget{U+004E}{U+004E}
&N&Latin Capital Letter N&\oldsquare&oldsquare&
\hyperlink{U+25A1}{U+25A1}
&
\\
\hypertarget{U+004F}{U+004F}
&O&Latin Capital Letter O&\pentagramroughinvsolid&pentagramroughinvsolid&
&
\hyperlink{U+006F}{U+006F}
\\
\hypertarget{U+0050}{U+0050}
&P&Latin Capital Letter P&\pentagramintRtriangle&pentagramintRtriangle&
&
\\
\hypertarget{U+0051}{U+0051}
&Q&Latin Capital Letter Q&\pentagraminvcircled&pentagraminvcircled&
&
\hyperlink{U+0071}{U+0071}
\\
\hypertarget{U+0052}{U+0052}
&R&Latin Capital Letter R&\pentagraminvintL&pentagraminvintL&
&
\hyperlink{U+0154}{U+0154}
\\
\hypertarget{U+0053}{U+0053}
&S&Latin Capital Letter S&\whiteknifenorth&whiteknifenorth&
&
\\
\hypertarget{U+0054}{U+0054}
&T&Latin Capital Letter T&\pentagramroughinvcircled&pentagramroughinvcircled&
&
\\
\hypertarget{U+0055}{U+0055}
&U&Latin Capital Letter U&\oldpentagraminvsolid&oldpentagraminvsolid&
\hyperlink{U+00CD}{U+00CD}
&
\\
\hypertarget{U+0056}{U+0056}
&V&Latin Capital Letter V&\oldopposition&oldopposition&
\hyperlink{U+260D}{U+260D}
&
\\
\hypertarget{U+0057}{U+0057}
&W&Latin Capital Letter W&\oldpentagraminv&oldpentagraminv&
\hyperlink{U+26E7}{U+26E7}
&
\\
\hypertarget{U+0058}{U+0058}
&X&Latin Capital Letter X&\oldascendingnode&oldascendingnode&
\hyperlink{U+260A}{U+260A}
&
\\
\hypertarget{U+0059}{U+0059}
&Y&Latin Capital Letter Y&\pentagramroughinv&pentagramroughinv&
&
\\
\hypertarget{U+005A}{U+005A}
&Z&Latin Capital Letter Z&\olddescendingnode&olddescendingnode&
\hyperlink{U+260B}{U+260B}
&
\\
\hypertarget{U+005B}{U+005B}
&[&Left Square Bracket&\fireelement&fireelement&
&
\\
\hypertarget{U+005C}{U+005C}
&\textbackslash&Reverse Solidus&\serpentimpaled&serpentimpaled&
&
\\
\hypertarget{U+005D}{U+005D}
&]&Right Square Bracket&\airelement&airelement&
&
\\
\hypertarget{U+005E}{U+005E}
&\^{}&Circumflex Accent&\moonlastquarteroutline&moonlastquarteroutline&
&
\\
\hypertarget{U+005F}{U+005F}
&\_&Low Line&\moonlastquartersolid&moonlastquartersolid&
&
\\
\hypertarget{U+0060}{U+0060}
&\`{}&Grave Accent&\oldmercury&oldmercury&
\hyperlink{U+263F}{U+263F}
&
\\
\hypertarget{U+0061}{U+0061}
&a&Latin Small Letter A&\athamesouth&athamesouth&
&
\\
\hypertarget{U+0062}{U+0062}
&b&Latin Small Letter B&\olduranus&olduranus&
\hyperlink{U+2645}{U+2645}
&
\\
\hypertarget{U+0063}{U+0063}
&c&Latin Small Letter C&\oldjupiter&oldjupiter&
\hyperlink{U+2643}{U+2643}
&
\\
\hypertarget{U+0064}{U+0064}
&d&Latin Small Letter D&\athameeast&athameeast&
&
\\
\hypertarget{U+0065}{U+0065}
&e&Latin Small Letter E&\pentagramintRcircled&pentagramintRcircled&
&
\hyperlink{U+00E9}{U+00E9}
\\
\hypertarget{U+0066}{U+0066}
&f&Latin Small Letter F&\athamewest&athamewest&
&
\\
\hypertarget{U+0067}{U+0067}
&g&Latin Small Letter G&\hammer&hammer&
&
\\
\hypertarget{U+0068}{U+0068}
&h&Latin Small Letter H&\awen&awen&
&
\\
\end{tabular}
\par
\begin{tabular}{@{}
  >{\ttfamily}l@{}
  >{\sffamily}c@{\hspace{2mm}}
  >{\sffamily\raggedright\vstrut}p{.3\textwidth}
  >{\Large\symbats}c
  >{\ttfamily\raggedright\arraybackslash\textbackslash}m{.29\textwidth}
  >{\ttfamily}l@{\hspace{2mm}}
  >{\ttfamily}l@{}}
\normalfont\bfseries Code&
\normalfont\bfseries Char&
\normalfont\bfseries Unicode name&
\normalfont\normalsize\bfseries Glyph&
\multicolumn1{>{\normalfont\bfseries}l}{\LaTeX\ name}&
\normalfont\bfseries Dupes&
\normalfont\bfseries Compl\\\hline\vrule height1.75em width0pt
\hypertarget{U+0069}{U+0069}
&i&Latin Small Letter I&\oldpentagramintRsolid&oldpentagramintRsolid&
\hyperlink{U+00ED}{U+00ED}
&
\\
\hypertarget{U+006A}{U+006A}
&j&Latin Small Letter J&\fylfotilkley&fylfotilkley&
&
\\
\hypertarget{U+006B}{U+006B}
&k&Latin Small Letter K&\aegishjalmur&aegishjalmur&
&
\\
\hypertarget{U+006C}{U+006C}
&l&Latin Small Letter L&\eightsabbats&eightsabbats&
&
\\
\hypertarget{U+006D}{U+006D}
&m&Latin Small Letter M&\oldpluto&oldpluto&
\hyperlink{U+2647}{U+2647}
&
\\
\hypertarget{U+006E}{U+006E}
&n&Latin Small Letter N&\oldneptune&oldneptune&
\hyperlink{U+2646}{U+2646}
&
\\
\hypertarget{U+006F}{U+006F}
&o&Latin Small Letter O&\pentagramroughsolid&pentagramroughsolid&
&
\hyperlink{U+004F}{U+004F}
\\
\hypertarget{U+0070}{U+0070}
&p&Latin Small Letter P&\pentagramwithtriangle&pentagramwithtriangle&
&
\\
\hypertarget{U+0071}{U+0071}
&q&Latin Small Letter Q&\pentagramcircled&pentagramcircled&
&
\hyperlink{U+0051}{U+0051}
\\
\hypertarget{U+0072}{U+0072}
&r&Latin Small Letter R&\oldpentagramintR&oldpentagramintR&
\hyperlink{U+26E5}{U+26E5}
&
\hyperlink{U+0155}{U+0155}
\\
\hypertarget{U+0073}{U+0073}
&s&Latin Small Letter S&\athamenorth&athamenorth&
&
\\
\hypertarget{U+0074}{U+0074}
&t&Latin Small Letter T&\pentagramroughcircled&pentagramroughcircled&
&
\\
\hypertarget{U+0075}{U+0075}
&u&Latin Small Letter U&\pentagramsolid&pentagramsolid&
&
\\
\hypertarget{U+0076}{U+0076}
&v&Latin Small Letter V&\oldsaturn&oldsaturn&
\hyperlink{U+2644}{U+2644}
&
\\
\hypertarget{U+0077}{U+0077}
&w&Latin Small Letter W&\oldpentagram&oldpentagram&
\hyperlink{U+26E4}{U+26E4}
&
\\
\hypertarget{U+0078}{U+0078}
&x&Latin Small Letter X&\mars&mars&
&
\\
\hypertarget{U+0079}{U+0079}
&y&Latin Small Letter Y&\pentagramrough&pentagramrough&
&
\\
\hypertarget{U+007A}{U+007A}
&z&Latin Small Letter Z&\venus&venus&
&
\\
\hypertarget{U+007B}{U+007B}
&\{&Left Curly Bracket&\waterelement&waterelement&
&
\\
\hypertarget{U+007C}{U+007C}
&\|&Vertical Line&\celtictrefoilknotsolid&celtictrefoilknotsolid&
&
\\
\hypertarget{U+007D}{U+007D}
&\}&Right Curly Bracket&\earthelement&earthelement&
&
\\
\hypertarget{U+007E}{U+007E}
&\~{}&Tilde&\vargemini&vargemini&
&
\\
\hypertarget{U+00A1}{U+00A1}
&¡&Inv Exclamation Mark&\oldaries&oldaries&
\hyperlink{U+2648}{U+2648}
&
\\
\hypertarget{U+00A2}{U+00A2}
&¢&Cent Sign&\oldcancer&oldcancer&
\hyperlink{U+264B}{U+264B}
&
\\
\hypertarget{U+00A3}{U+00A3}
&£&Pound Sign&\oldmoonwaxingoutline&oldmoonwaxingoutline&
\hyperlink{U+263D}{U+263D}
&
\\
\hypertarget{U+00A7}{U+00A7}
&§&Section Sign&\oldvirgo&oldvirgo&
\hyperlink{U+264D}{U+264D}
&
\\
\hypertarget{U+00AA}{U+00AA}
&ª&Feminine Ordinal Indicator&\oldsaggitarius&oldsaggitarius&
\hyperlink{U+2650}{U+2650}
&
\\
\hypertarget{U+00B6}{U+00B6}
&¶&Pilcrow Sign&\oldlibra&oldlibra&
\hyperlink{U+264E}{U+264E}
&
\\
\hypertarget{U+00BA}{U+00BA}
&º&Masculine Ordinal Indicator&\oldcapricorn&oldcapricorn&
\hyperlink{U+2651}{U+2651}
&
\\
\hypertarget{U+00C0}{U+00C0}
&À&Latin Capital Letter A with Grave&\powergoingforth&powergoingforth&
&
\\
\hypertarget{U+00C5}{U+00C5}
&Å&Latin Capital Letter A with Ring Above&\samhuinn&samhuinn&
&
\\
\hypertarget{U+00C8}{U+00C8}
&È&Latin Capital Letter E with Grave&\parting&parting&
&
\\
\hypertarget{U+00C9}{U+00C9}
&É&Latin Capital Letter E with Acute&\pentagraminvintRcircled&pentagraminvintRcircled&
&
\hyperlink{U+0045}{U+0045}
\\
\hypertarget{U+00CC}{U+00CC}
&Ì&Latin Capital Letter I with Grave&\beltane&beltane&
&
\\
\hypertarget{U+00CD}{U+00CD}
&Í&Latin Capital Letter I with Acute&\pentagraminvsolid&pentagraminvsolid&
\hyperlink{U+0055}{U+0055}
&
\\
\hypertarget{U+00CE}{U+00CE}
&Î&Latin Capital Letter I with Circumflex&\imbolc&imbolc&
&
\\
\end{tabular}
\par
\begin{tabular}{@{}
  >{\ttfamily}l@{}
  >{\sffamily}c@{\hspace{2mm}}
  >{\sffamily\raggedright\vstrut}p{.3\textwidth}
  >{\Large\symbats}c
  >{\ttfamily\raggedright\arraybackslash\textbackslash}m{.29\textwidth}
  >{\ttfamily}l@{\hspace{2mm}}
  >{\ttfamily}l@{}}
\normalfont\bfseries Code&
\normalfont\bfseries Char&
\normalfont\bfseries Unicode name&
\normalfont\normalsize\bfseries Glyph&
\multicolumn1{>{\normalfont\bfseries}l}{\LaTeX\ name}&
\normalfont\bfseries Dupes&
\normalfont\bfseries Compl\\\hline\vrule height1.75em width0pt
\hypertarget{U+00CF}{U+00CF}
&Ï&Latin Capital Letter I with Diaeresis&\eostre&eostre&
&
\\
\hypertarget{U+00D3}{U+00D3}
&Ó&Latin Capital Letter O with Acute&\venusfigureoutline&venusfigureoutline&
&
\\
\hypertarget{U+00D4}{U+00D4}
&Ô&Latin Capital Letter O with Circumflex&\lughnasadh&lughnasadh&
&
\\
\hypertarget{U+00D9}{U+00D9}
&Ù&Latin Capital Letter U with Grave&\perfectcouplethirddegree&perfectcouplethirddegree&
&
\\
\hypertarget{U+00E0}{U+00E0}
&à&Latin Small Letter A with Grave&\oldankh&oldankh&
\hyperlink{U+2625}{U+2625}
&
\\
\hypertarget{U+00E8}{U+00E8}
&è&Latin Small Letter E with Grave&\salute&salute&
&
\\
\hypertarget{U+00E9}{U+00E9}
&é&Latin Small Letter E with Acute&\pentagramintLcircled&pentagramintLcircled&
&
\hyperlink{U+0065}{U+0065}
\\
\hypertarget{U+00EC}{U+00EC}
&ì&Latin Small Letter I with Grave&\scourge&scourge&
&
\\
\hypertarget{U+00ED}{U+00ED}
&í&Latin Small Letter I with Acute&\pentagramintRsolid&pentagramintRsolid&
\hyperlink{U+0069}{U+0069}
&
\\
\hypertarget{U+00F2}{U+00F2}
&ò&Latin Small Letter O with Grave&\goddessasmoon&goddessasmoon&
&
\\
\hypertarget{U+00F3}{U+00F3}
&ó&Latin Small Letter O with Acute&\venusfiguresolid&venusfiguresolid&
&
\\
\hypertarget{U+00F9}{U+00F9}
&ù&Latin Small Letter U with Grave&\perfectcouple&perfectcouple&
&
\\
\hypertarget{U+0132}{U+0132}
&IJ&Latin Capital Ligature Ij&\yule&yule&
&
\\
\hypertarget{U+0154}{U+0154}
&Ŕ&Latin Capital Letter R with Acute&\pentagraminvintR&pentagraminvintR&
&
\hyperlink{U+0052}{U+0052}
\\
\hypertarget{U+0155}{U+0155}
&ŕ&Latin Small Letter R with Acute&\oldpentagramintL&oldpentagramintL&
\hyperlink{U+26E6}{U+26E6}
&
\\
\hypertarget{U+1681}{U+1681}
&ᚁ&Ogham Letter Beith&\OGbeith&OGbeith&
&
\\
\hypertarget{U+1682}{U+1682}
&ᚂ&Ogham Letter Luis&\OGluis&OGluis&
&
\\
\hypertarget{U+1683}{U+1683}
&ᚃ&Ogham Letter Fearn&\OGfearn&OGfearn&
&
\\
\hypertarget{U+1684}{U+1684}
&ᚄ&Ogham Letter Sail&\OGsail&OGsail&
&
\\
\hypertarget{U+1685}{U+1685}
&ᚅ&Ogham Letter Nion&\OGnion&OGnion&
&
\\
\hypertarget{U+1686}{U+1686}
&ᚆ&Ogham Letter Uath&\OGuath&OGuath&
&
\\
\hypertarget{U+1687}{U+1687}
&ᚇ&Ogham Letter Dair&\OGdair&OGdair&
&
\\
\hypertarget{U+1688}{U+1688}
&ᚈ&Ogham Letter Tinne&\OGtinne&OGtinne&
&
\\
\hypertarget{U+1689}{U+1689}
&ᚉ&Ogham Letter Coll&\OGcoll&OGcoll&
&
\\
\hypertarget{U+168A}{U+168A}
&ᚊ&Ogham Letter Ceirt&\OGceirt&OGceirt&
&
\\
\hypertarget{U+168B}{U+168B}
&ᚋ&Ogham Letter Muin&\OGmuin&OGmuin&
&
\\
\hypertarget{U+168C}{U+168C}
&ᚌ&Ogham Letter Gort&\OGgort&OGgort&
&
\\
\hypertarget{U+168D}{U+168D}
&ᚍ&Ogham Letter Ngeadal&\OGngeadal&OGngeadal&
&
\\
\hypertarget{U+168E}{U+168E}
&ᚎ&Ogham Letter Straif&\OGstraif&OGstraif&
&
\\
\hypertarget{U+168F}{U+168F}
&ᚏ&Ogham Letter Ruis&\OGruis&OGruis&
&
\\
\hypertarget{U+1690}{U+1690}
&ᚐ&Ogham Letter Ailm&\OGailm&OGailm&
&
\\
\hypertarget{U+1691}{U+1691}
&ᚑ&Ogham Letter Onn&\OGonn&OGonn&
&
\\
\hypertarget{U+1692}{U+1692}
&ᚒ&Ogham Letter Ur&\OGur&OGur&
&
\\
\hypertarget{U+1693}{U+1693}
&ᚓ&Ogham Letter Eadhadh&\OGeadhadh&OGeadhadh&
&
\\
\hypertarget{U+1694}{U+1694}
&ᚔ&Ogham Letter Iodhadh&\OGiodhadh&OGiodhadh&
&
\\
\hypertarget{U+1695}{U+1695}
&ᚕ&Ogham Letter Eabhadh&\OGeabhadh&OGeabhadh&
&
\\
\end{tabular}
\par
\begin{tabular}{@{}
  >{\ttfamily}l@{}
  >{\sffamily}c@{\hspace{2mm}}
  >{\sffamily\raggedright\vstrut}p{.3\textwidth}
  >{\Large\symbats}c
  >{\ttfamily\raggedright\arraybackslash\textbackslash}m{.29\textwidth}
  >{\ttfamily}l@{\hspace{2mm}}
  >{\ttfamily}l@{}}
\normalfont\bfseries Code&
\normalfont\bfseries Char&
\normalfont\bfseries Unicode name&
\normalfont\normalsize\bfseries Glyph&
\multicolumn1{>{\normalfont\bfseries}l}{\LaTeX\ name}&
\normalfont\bfseries Dupes&
\normalfont\bfseries Compl\\\hline\vrule height1.75em width0pt
\hypertarget{U+1696}{U+1696}
&ᚖ&Ogham Letter Or&\OGor&OGor&
&
\\
\hypertarget{U+1697}{U+1697}
&ᚗ&Ogham Letter Uilleann&\OGuilleann&OGuilleann&
&
\\
\hypertarget{U+1698}{U+1698}
&ᚘ&Ogham Letter Ifin&\OGifin&OGifin&
&
\\
\hypertarget{U+1699}{U+1699}
&ᚙ&Ogham Letter Eamhancholl&\OGeamhancholl&OGeamhancholl&
&
\\
\hypertarget{U+169A}{U+169A}
&ᚚ&Ogham Letter Peith&\OGpeith&OGpeith&
&
\\
\hypertarget{U+169B}{U+169B}
&᚛&Ogham Feather Mark&\OGfeather&OGfeather&
&
\\
\hypertarget{U+169C}{U+169C}
&᚜&Ogham Reversed Feather Mark&\OGreversedfeather&OGreversedfeather&
&
\\
\hypertarget{U+16A0}{U+16A0}
&ᚠ&Runic Letter Fehu Feoh Fe F&\RUfehufeohfef&RUfehufeohfef&
&
\\
\hypertarget{U+16A1}{U+16A1}
&ᚡ&Runic Letter V&\RUv&RUv&
&
\\
\hypertarget{U+16A2}{U+16A2}
&ᚢ&Runic Letter Uruz Ur U&\RUuruzuru&RUuruzuru&
&
\\
\hypertarget{U+16A3}{U+16A3}
&ᚣ&Runic Letter Yr&\RUyr&RUyr&
&
\\
\hypertarget{U+16A4}{U+16A4}
&ᚤ&Runic Letter Y&\RUy&RUy&
&
\\
\hypertarget{U+16A5}{U+16A5}
&ᚥ&Runic Letter W&\RUw&RUw&
&
\\
\hypertarget{U+16A6}{U+16A6}
&ᚦ&Runic Letter Thurisaz Thurs Thorn&\RUthurisazthursthorn&RUthurisazthursthorn&
&
\\
\hypertarget{U+16A7}{U+16A7}
&ᚧ&Runic Letter Eth&\RUeth&RUeth&
&
\\
\hypertarget{U+16A8}{U+16A8}
&ᚨ&Runic Letter Ansuz A&\RUansuza&RUansuza&
&
\\
\hypertarget{U+16A9}{U+16A9}
&ᚩ&Runic Letter Os O&\RUoso&RUoso&
&
\\
\hypertarget{U+16AA}{U+16AA}
&ᚪ&Runic Letter Ac A&\RUaca&RUaca&
&
\\
\hypertarget{U+16AB}{U+16AB}
&ᚫ&Runic Letter Aesc&\RUaesc&RUaesc&
&
\\
\hypertarget{U+16AC}{U+16AC}
&ᚬ&Runic Letter Long-Branch-Oss O&\RUlongbranchosso&RUlongbranchosso&
&
\\
\hypertarget{U+16AD}{U+16AD}
&ᚭ&Runic Letter Short-Twig-Oss O&\RUshorttwigosso&RUshorttwigosso&
&
\\
\hypertarget{U+16AE}{U+16AE}
&ᚮ&Runic Letter O&\RUo&RUo&
&
\\
\hypertarget{U+16AF}{U+16AF}
&ᚯ&Runic Letter Oe&\RUoe&RUoe&
&
\\
\hypertarget{U+16B0}{U+16B0}
&ᚰ&Runic Letter On&\RUon&RUon&
&
\\
\hypertarget{U+16B1}{U+16B1}
&ᚱ&Runic Letter Raido Rad Reid R&\RUraidoradreidr&RUraidoradreidr&
&
\\
\hypertarget{U+16B2}{U+16B2}
&ᚲ&Runic Letter Kauna&\RUkauna&RUkauna&
&
\\
\hypertarget{U+16B3}{U+16B3}
&ᚳ&Runic Letter Cen&\RUcen&RUcen&
&
\\
\hypertarget{U+16B4}{U+16B4}
&ᚴ&Runic Letter Kaun K&\RUkaunk&RUkaunk&
&
\\
\hypertarget{U+16B5}{U+16B5}
&ᚵ&Runic Letter G&\RUg&RUg&
&
\\
\hypertarget{U+16B6}{U+16B6}
&ᚶ&Runic Letter Eng&\RUeng&RUeng&
&
\\
\hypertarget{U+16B7}{U+16B7}
&ᚷ&Runic Letter Gebo Gyfu G&\RUgebogyfug&RUgebogyfug&
&
\\
\hypertarget{U+16B8}{U+16B8}
&ᚸ&Runic Letter Gar&\RUgar&RUgar&
&
\\
\hypertarget{U+16B9}{U+16B9}
&ᚹ&Runic Letter Wunjo Wynn W&\RUwunjowynnw&RUwunjowynnw&
&
\\
\hypertarget{U+16BA}{U+16BA}
&ᚺ&Runic Letter Haglaz H&\RUhaglazh&RUhaglazh&
&
\\
\hypertarget{U+16BB}{U+16BB}
&ᚻ&Runic Letter Haegl H&\RUhaeglh&RUhaeglh&
&
\\
\hypertarget{U+16BC}{U+16BC}
&ᚼ&Runic Letter Long-Branch-Hagall H&\RUlongbranchhagallh&RUlongbranchhagallh&
&
\\
\end{tabular}
\par
\begin{tabular}{@{}
  >{\ttfamily}l@{}
  >{\sffamily}c@{\hspace{2mm}}
  >{\sffamily\raggedright\vstrut}p{.3\textwidth}
  >{\Large\symbats}c
  >{\ttfamily\raggedright\arraybackslash\textbackslash}m{.29\textwidth}
  >{\ttfamily}l@{\hspace{2mm}}
  >{\ttfamily}l@{}}
\normalfont\bfseries Code&
\normalfont\bfseries Char&
\normalfont\bfseries Unicode name&
\normalfont\normalsize\bfseries Glyph&
\multicolumn1{>{\normalfont\bfseries}l}{\LaTeX\ name}&
\normalfont\bfseries Dupes&
\normalfont\bfseries Compl\\\hline\vrule height1.75em width0pt
\hypertarget{U+16BD}{U+16BD}
&ᚽ&Runic Letter Short-Twig-Hagall H&\RUshorttwighagallh&RUshorttwighagallh&
&
\\
\hypertarget{U+16BE}{U+16BE}
&ᚾ&Runic Letter Naudiz Nyd Naud N&\RUnaudiznydnaudn&RUnaudiznydnaudn&
&
\\
\hypertarget{U+16BF}{U+16BF}
&ᚿ&Runic Letter Short-Twig-Naud N&\RUshorttwignaudn&RUshorttwignaudn&
&
\\
\hypertarget{U+16C0}{U+16C0}
&ᛀ&Runic Letter Dotted-N&\RUdottedn&RUdottedn&
&
\\
\hypertarget{U+16C1}{U+16C1}
&ᛁ&Runic Letter Isaz Is Iss I&\RUisazisissi&RUisazisissi&
&
\\
\hypertarget{U+16C2}{U+16C2}
&ᛂ&Runic Letter E&\RUe&RUe&
&
\\
\hypertarget{U+16C3}{U+16C3}
&ᛃ&Runic Letter Jeran J&\RUjeranj&RUjeranj&
&
\\
\hypertarget{U+16C4}{U+16C4}
&ᛄ&Runic Letter Ger&\RUger&RUger&
&
\\
\hypertarget{U+16C5}{U+16C5}
&ᛅ&Runic Letter Long-Branch-Ar Ae&\RUlongbrancharae&RUlongbrancharae&
&
\\
\hypertarget{U+16C6}{U+16C6}
&ᛆ&Runic Letter Short-Twig-Ar A&\RUshorttwigara&RUshorttwigara&
&
\\
\hypertarget{U+16C7}{U+16C7}
&ᛇ&Runic Letter Iwaz Eoh&\RUiwazeoh&RUiwazeoh&
&
\\
\hypertarget{U+16C8}{U+16C8}
&ᛈ&Runic Letter Pertho Peorth P&\RUperthopeorthp&RUperthopeorthp&
&
\\
\hypertarget{U+16C9}{U+16C9}
&ᛉ&Runic Letter Algiz Eolhx&\RUalgizeolhx&RUalgizeolhx&
&
\\
\hypertarget{U+16CA}{U+16CA}
&ᛊ&Runic Letter Sowilo S&\RUsowilos&RUsowilos&
&
\\
\hypertarget{U+16CB}{U+16CB}
&ᛋ&Runic Letter Sigel Long-Branch-Sol S&\RUsigellongbranchsols&RUsigellongbranchsols&
&
\\
\hypertarget{U+16CC}{U+16CC}
&ᛌ&Runic Letter Short-Twig-Sol S&\RUshorttwigsols&RUshorttwigsols&
&
\\
\hypertarget{U+16CD}{U+16CD}
&ᛍ&Runic Letter C&\RUc&RUc&
&
\\
\hypertarget{U+16CE}{U+16CE}
&ᛎ&Runic Letter Z&\RUz&RUz&
&
\\
\hypertarget{U+16CF}{U+16CF}
&ᛏ&Runic Letter Tiwaz Tir Tyr T&\RUtiwaztirtyrt&RUtiwaztirtyrt&
&
\\
\hypertarget{U+16D0}{U+16D0}
&ᛐ&Runic Letter Short-Twig-Tyr T&\RUshorttwigtyrt&RUshorttwigtyrt&
&
\\
\hypertarget{U+16D1}{U+16D1}
&ᛑ&Runic Letter D&\RUd&RUd&
&
\\
\hypertarget{U+16D2}{U+16D2}
&ᛒ&Runic Letter Berkanan Beorc Bjarkan B&\RUberkananbeorcbjarkanb&RUberkananbeorcbjarkanb&
&
\\
\hypertarget{U+16D3}{U+16D3}
&ᛓ&Runic Letter Short-Twig-Bjarkan B&\RUshorttwigbjarkanb&RUshorttwigbjarkanb&
&
\\
\hypertarget{U+16D4}{U+16D4}
&ᛔ&Runic Letter Dotted-P&\RUdottedp&RUdottedp&
&
\\
\hypertarget{U+16D5}{U+16D5}
&ᛕ&Runic Letter Open-P&\RUopenp&RUopenp&
&
\\
\hypertarget{U+16D6}{U+16D6}
&ᛖ&Runic Letter Ehwaz Eh E&\RUehwazehe&RUehwazehe&
&
\\
\hypertarget{U+16D7}{U+16D7}
&ᛗ&Runic Letter Mannaz Man M&\RUmannazmanm&RUmannazmanm&
&
\\
\hypertarget{U+16D8}{U+16D8}
&ᛘ&Runic Letter Long-Branch-Madr M&\RUlongbranchmadrm&RUlongbranchmadrm&
&
\\
\hypertarget{U+16D9}{U+16D9}
&ᛙ&Runic Letter Short-Twig-Madr M&\RUshorttwigmadrm&RUshorttwigmadrm&
&
\\
\hypertarget{U+16DA}{U+16DA}
&ᛚ&Runic Letter Laukaz Lagu Logr L&\RUlaukazlagulogrl&RUlaukazlagulogrl&
&
\\
\hypertarget{U+16DB}{U+16DB}
&ᛛ&Runic Letter Dotted-L&\RUdottedl&RUdottedl&
&
\\
\hypertarget{U+16DC}{U+16DC}
&ᛜ&Runic Letter Ingwaz&\RUingwaz&RUingwaz&
&
\\
\hypertarget{U+16DD}{U+16DD}
&ᛝ&Runic Letter Ing&\RUing&RUing&
&
\\
\hypertarget{U+16DE}{U+16DE}
&ᛞ&Runic Letter Dagaz Daeg D&\RUdagazdaegd&RUdagazdaegd&
&
\\
\hypertarget{U+16DF}{U+16DF}
&ᛟ&Runic Letter Othalan Ethel O&\RUothalanethelo&RUothalanethelo&
&
\\
\hypertarget{U+16E0}{U+16E0}
&ᛠ&Runic Letter Ear&\RUear&RUear&
&
\\
\end{tabular}
\par
\begin{tabular}{@{}
  >{\ttfamily}l@{}
  >{\sffamily}c@{\hspace{2mm}}
  >{\sffamily\raggedright\vstrut}p{.3\textwidth}
  >{\Large\symbats}c
  >{\ttfamily\raggedright\arraybackslash\textbackslash}m{.29\textwidth}
  >{\ttfamily}l@{\hspace{2mm}}
  >{\ttfamily}l@{}}
\normalfont\bfseries Code&
\normalfont\bfseries Char&
\normalfont\bfseries Unicode name&
\normalfont\normalsize\bfseries Glyph&
\multicolumn1{>{\normalfont\bfseries}l}{\LaTeX\ name}&
\normalfont\bfseries Dupes&
\normalfont\bfseries Compl\\\hline\vrule height1.75em width0pt
\hypertarget{U+16E1}{U+16E1}
&ᛡ&Runic Letter Ior&\RUior&RUior&
&
\\
\hypertarget{U+16E2}{U+16E2}
&ᛢ&Runic Letter Cweorth&\RUcweorth&RUcweorth&
&
\\
\hypertarget{U+16E3}{U+16E3}
&ᛣ&Runic Letter Calc&\RUcalc&RUcalc&
&
\\
\hypertarget{U+16E4}{U+16E4}
&ᛤ&Runic Letter Cealc&\RUcealc&RUcealc&
&
\\
\hypertarget{U+16E5}{U+16E5}
&ᛥ&Runic Letter Stan&\RUstan&RUstan&
&
\\
\hypertarget{U+16E6}{U+16E6}
&ᛦ&Runic Letter Long-Branch-Yr&\RUlongbranchyr&RUlongbranchyr&
&
\\
\hypertarget{U+16E7}{U+16E7}
&ᛧ&Runic Letter Short-Twig-Yr&\RUshorttwigyr&RUshorttwigyr&
&
\\
\hypertarget{U+16E8}{U+16E8}
&ᛨ&Runic Letter Icelandic-Yr&\RUicelandicyr&RUicelandicyr&
&
\\
\hypertarget{U+16E9}{U+16E9}
&ᛩ&Runic Letter Q&\RUq&RUq&
&
\\
\hypertarget{U+16EA}{U+16EA}
&ᛪ&Runic Letter X&\RUx&RUx&
&
\\
\hypertarget{U+16EB}{U+16EB}
&᛫&Runic Single Punctuation&\RUsinglepunctuation&RUsinglepunctuation&
&
\\
\hypertarget{U+16EC}{U+16EC}
&᛬&Runic Multiple Punctuation&\RUmultiplepunctuation&RUmultiplepunctuation&
&
\\
\hypertarget{U+16ED}{U+16ED}
&᛭&Runic Cross Punctuation&\RUcrosspunctuation&RUcrosspunctuation&
&
\\
\hypertarget{U+16EE}{U+16EE}
&ᛮ&Runic Arlaug Symbol&\RUarlaugsymbol&RUarlaugsymbol&
&
\\
\hypertarget{U+16EF}{U+16EF}
&ᛯ&Runic Tvimadur Symbol&\RUtvimadursymbol&RUtvimadursymbol&
&
\\
\hypertarget{U+16F0}{U+16F0}
&ᛰ&Runic Belgthor Symbol&\RUbelgthorsymbol&RUbelgthorsymbol&
&
\\
\hypertarget{U+16F1}{U+16F1}
&ᛱ&Runic Letter K&\RUk&RUk&
&
\\
\hypertarget{U+16F2}{U+16F2}
&ᛲ&Runic Letter Sh&\RUsh&RUsh&
&
\\
\hypertarget{U+16F3}{U+16F3}
&ᛳ&Runic Letter Oo&\RUoo&RUoo&
&
\\
\hypertarget{U+16F4}{U+16F4}
&ᛴ&Runic Letter Franks Casket Os&\RUfrankscasketos&RUfrankscasketos&
&
\\
\hypertarget{U+16F5}{U+16F5}
&ᛵ&Runic Letter Franks Casket Is&\RUfrankscasketis&RUfrankscasketis&
&
\\
\hypertarget{U+16F6}{U+16F6}
&ᛶ&Runic Letter Franks Casket Eh&\RUfrankscasketeh&RUfrankscasketeh&
&
\\
\hypertarget{U+16F7}{U+16F7}
&ᛷ&Runic Letter Franks Casket Ac&\RUfrankscasketac&RUfrankscasketac&
&
\\
\hypertarget{U+16F8}{U+16F8}
&ᛸ&Runic Letter Franks Casket Aesc&\RUfrankscasketaesc&RUfrankscasketaesc&
&
\\
\hypertarget{U+2013}{U+2013}
&–&En Dash&\oldaquarius&oldaquarius&
\hyperlink{U+2652}{U+2652}
&
\\
\hypertarget{U+2022}{U+2022}
&•&Bullet&\scorpio&scorpio&
&
\\
\hypertarget{U+211E}{U+211E}
&℞&Prescription Take (Recipe)&\prescription&prescription&
&
\\
\hypertarget{U+2122}{U+2122}
&™&Trade Mark Sign&\oldtaurus&oldtaurus&
\hyperlink{U+2649}{U+2649}
&
\\
\hypertarget{U+221E}{U+221E}
&∞&Infinity&\oldleo&oldleo&
\hyperlink{U+264C}{U+264C}
&
\\
\hypertarget{U+2220}{U+2220}
&∠&Angle&\angleacute&angleacute&
&
\\
\hypertarget{U+2260}{U+2260}
&≠&Not Equal To&\oldpisces&oldpisces&
\hyperlink{U+2653}{U+2653}
&
\\
\hypertarget{U+2295}{U+2295}
&⊕&Circled Plus&\circleplus&circleplus&
&
\\
\hypertarget{U+25A1}{U+25A1}
&□&White Square&\square&square&
\hyperlink{U+004F}{U+004F}
&
\\
\hypertarget{U+2609}{U+2609}
&☉&Sun&\dupsun&dupsun&
&
\\
\hypertarget{U+260A}{U+260A}
&☊&Ascending Node&\ascendingnode&ascendingnode&
\hyperlink{U+0058}{U+0058}
&
\\
\hypertarget{U+260B}{U+260B}
&☋&Descending Node&\descendingnode&descendingnode&
\hyperlink{U+005A}{U+005A}
&
\\
\end{tabular}
\par
\begin{tabular}{@{}
  >{\ttfamily}l@{}
  >{\sffamily}c@{\hspace{2mm}}
  >{\sffamily\raggedright\vstrut}p{.3\textwidth}
  >{\Large\symbats}c
  >{\ttfamily\raggedright\arraybackslash\textbackslash}m{.29\textwidth}
  >{\ttfamily}l@{\hspace{2mm}}
  >{\ttfamily}l@{}}
\normalfont\bfseries Code&
\normalfont\bfseries Char&
\normalfont\bfseries Unicode name&
\normalfont\normalsize\bfseries Glyph&
\multicolumn1{>{\normalfont\bfseries}l}{\LaTeX\ name}&
\normalfont\bfseries Dupes&
\normalfont\bfseries Compl\\\hline\vrule height1.75em width0pt
\hypertarget{U+260C}{U+260C}
&☌&Conjunction&\conjunction&conjunction&
\hyperlink{U+0043}{U+0043}
&
\\
\hypertarget{U+260D}{U+260D}
&☍&Opposition&\opposition&opposition&
\hyperlink{U+0056}{U+0056}
&
\\
\hypertarget{U+2625}{U+2625}
&☥&Ankh&\ankh&ankh&
\hyperlink{U+00E0}{U+00E0}
&
\\
\hypertarget{U+263C}{U+263C}
&☼&White Sun with Rays&\sunrays&sunrays&
&
\\
\hypertarget{U+263D}{U+263D}
&☽&First Quarter Moon&\moonwaxingoutline&moonwaxingoutline&
\hyperlink{U+00A3}{U+00A3}
&
\\
\hypertarget{U+263E}{U+263E}
&☾&Last Quarter Moon&\moonwaningoutline&moonwaningoutline&
\hyperlink{U+0026}{U+0026}
&
\\
\hypertarget{U+263F}{U+263F}
&☿&Mercury&\mercury&mercury&
\hyperlink{U+0060}{U+0060}
&
\\
\hypertarget{U+2640}{U+2640}
&♀&Female Sign&\female&female&
&
\\
\hypertarget{U+2641}{U+2641}
&♁&Earth&\earth&earth&
\hyperlink{U+003C}{U+003C}
&
\\
\hypertarget{U+2642}{U+2642}
&♂&Male Sign&\male&male&
&
\\
\hypertarget{U+2643}{U+2643}
&♃&Jupiter&\jupiter&jupiter&
\hyperlink{U+0063}{U+0063}
&
\\
\hypertarget{U+2644}{U+2644}
&♄&Saturn&\saturn&saturn&
\hyperlink{U+0076}{U+0076}
&
\\
\hypertarget{U+2645}{U+2645}
&♅&Uranus&\uranus&uranus&
\hyperlink{U+0062}{U+0062}
&
\\
\hypertarget{U+2646}{U+2646}
&♆&Neptune&\neptune&neptune&
\hyperlink{U+006E}{U+006E}
&
\\
\hypertarget{U+2647}{U+2647}
&♇&Pluto&\pluto&pluto&
\hyperlink{U+006D}{U+006D}
&
\\
\hypertarget{U+2648}{U+2648}
&♈&Aries&\aries&aries&
\hyperlink{U+00A1}{U+00A1}
&
\\
\hypertarget{U+2649}{U+2649}
&♉&Taurus&\taurus&taurus&
\hyperlink{U+2122}{U+2122}
&
\\
\hypertarget{U+264A}{U+264A}
&♊&Gemini&\gemini&gemini&
\hyperlink{U+0023}{U+0023}
&
\\
\hypertarget{U+264B}{U+264B}
&♋&Cancer&\cancer&cancer&
\hyperlink{U+00A2}{U+00A2}
&
\\
\hypertarget{U+264C}{U+264C}
&♌&Leo&\leo&leo&
\hyperlink{U+221E}{U+221E}
&
\\
\hypertarget{U+264D}{U+264D}
&♍&Virgo&\virgo&virgo&
\hyperlink{U+00A7}{U+00A7}
&
\\
\hypertarget{U+264E}{U+264E}
&♎&Libra&\libra&libra&
\hyperlink{U+00B6}{U+00B6}
&
\\
\hypertarget{U+264F}{U+264F}
&♏&Scorpius&\scorpius&scorpius&
&
\\
\hypertarget{U+2650}{U+2650}
&♐&Sagittarius&\saggitarius&saggitarius&
\hyperlink{U+00AA}{U+00AA}
&
\\
\hypertarget{U+2651}{U+2651}
&♑&Capricorn&\capricorn&capricorn&
\hyperlink{U+00BA}{U+00BA}
&
\\
\hypertarget{U+2652}{U+2652}
&♒&Aquarius&\aquarius&aquarius&
\hyperlink{U+2013}{U+2013}
&
\\
\hypertarget{U+2653}{U+2653}
&♓&Pisces&\pisces&pisces&
\hyperlink{U+2260}{U+2260}
&
\\
\hypertarget{U+26A2}{U+26A2}
&⚢&Doubled Female Sign&\doubledfemale&doubledfemale&
&
\\
\hypertarget{U+26A3}{U+26A3}
&⚣&Doubled Male Sign&\doubledmale&doubledmale&
&
\\
\hypertarget{U+26A4}{U+26A4}
&⚤&Interlocked Female and Male Sign&\interlockedmalefemale&interlockedmalefemale&
&
\\
\hypertarget{U+26A5}{U+26A5}
&⚥&Male and Female Sign&\maleandfemale&maleandfemale&
&
\\
\hypertarget{U+26A6}{U+26A6}
&⚦&Male with Stroke Sign&\malewithstroke&malewithstroke&
&
\\
\hypertarget{U+26A7}{U+26A7}
&⚧&Male with Stroke and Male and Female Sign&\malewithstrokemalefemale&malewithstrokemalefemale&
&
\\
\hypertarget{U+26A8}{U+26A8}
&⚨&Vertical Male with Stroke Sign&\verticalmalewithstroke&verticalmalewithstroke&
&
\\
\hypertarget{U+26A9}{U+26A9}
&⚩&Horizontal Male with Stroke Sign&\horizontalmalewithstroke&horizontalmalewithstroke&
&
\\
\hypertarget{U+26B2}{U+26B2}
&⚲&Neuter&\neuter&neuter&
&
\\
\end{tabular}
\par
\begin{tabular}{@{}
  >{\ttfamily}l@{}
  >{\sffamily}c@{\hspace{2mm}}
  >{\sffamily\raggedright\vstrut}p{.3\textwidth}
  >{\Large\symbats}c
  >{\ttfamily\raggedright\arraybackslash\textbackslash}m{.29\textwidth}
  >{\ttfamily}l@{\hspace{2mm}}
  >{\ttfamily}l@{}}
\normalfont\bfseries Code&
\normalfont\bfseries Char&
\normalfont\bfseries Unicode name&
\normalfont\normalsize\bfseries Glyph&
\multicolumn1{>{\normalfont\bfseries}l}{\LaTeX\ name}&
\normalfont\bfseries Dupes&
\normalfont\bfseries Compl\\\hline\vrule height1.75em width0pt
\hypertarget{U+26B3}{U+26B3}
&⚳&Ceres&\ceres&ceres&
\hyperlink{U+002C}{U+002C}
&
\\
\hypertarget{U+26B4}{U+26B4}
&⚴&Pallas&\pallas&pallas&
&
\\
\hypertarget{U+26B5}{U+26B5}
&⚵&Juno&\juno&juno&
&
\\
\hypertarget{U+26B6}{U+26B6}
&⚶&Vesta&\vesta&vesta&
&
\\
\hypertarget{U+26B7}{U+26B7}
&⚷&Chiron&\chiron&chiron&
&
\\
\hypertarget{U+26B8}{U+26B8}
&⚸&Black Moon Lilith&\blackmoonlilith&blackmoonlilith&
&
\\
\hypertarget{U+26B9}{U+26B9}
&⚹&Sextile&\sextile&sextile&
\hyperlink{U+0042}{U+0042}
&
\\
\hypertarget{U+26BA}{U+26BA}
&⚺&Semisextile&\semisextile&semisextile&
&
\\
\hypertarget{U+26BB}{U+26BB}
&⚻&Quincunx&\quincunx&quincunx&
&
\\
\hypertarget{U+26BC}{U+26BC}
&⚼&Sesquiquadrate&\sesquiquadrate&sesquiquadrate&
&
\\
\hypertarget{U+26E2}{U+26E2}
&⛢&Astronomical Symbol For Uranus&\varuranus&varuranus&
&
\\
\hypertarget{U+26E4}{U+26E4}
&⛤&Pentagram&\pentagram&pentagram&
\hyperlink{U+0077}{U+0077}
&
\\
\hypertarget{U+26E5}{U+26E5}
&⛥&Right-Handed Interlaced Pentagram&\pentagramintR&pentagramintR&
&
\\
\hypertarget{U+26E6}{U+26E6}
&⛦&Left-Handed Interlaced Pentagram&\pentagramintL&pentagramintL&
\hyperlink{U+0155}{U+0155}
&
\\
\hypertarget{U+26E7}{U+26E7}
&⛧&Inv Pentagram&\pentagraminv&pentagraminv&
\hyperlink{U+0057}{U+0057}
&
\\
\hypertarget{U+2BC9}{U+2BC9}
&⯉&Neptune Form Two&\neptuneformtwo&neptuneformtwo&
&
\\
\hypertarget{U+2BD3}{U+2BD3}
&⯓&Pluto Form Two&\plutoformtwo&plutoformtwo&
&
\\
\hypertarget{U+2BD4}{U+2BD4}
&&Pluto Form Three&\plutoformthree&plutoformthree&
&
\\
\hypertarget{U+2BD5}{U+2BD5}
&⯕&Pluto Form Four&\plutoformfour&plutoformfour&
&
\\
\hypertarget{U+2BD6}{U+2BD6}
&⯖&Pluto Form Five&\plutoformfive&plutoformfive&
&
\\
\hypertarget{U+2BD7}{U+2BD7}
&⯗&Transpluto&\transpluto&transpluto&
&
\\
\hypertarget{U+2BD8}{U+2BD8}
&⯘&Proserpina&\proserpina&proserpina&
&
\\
\hypertarget{U+2BD9}{U+2BD9}
&⯙&Astraea&\astraea&astraea&
&
\\
\hypertarget{U+2BDA}{U+2BDA}
&⯚&Hygiea&\hygiea&hygiea&
&
\\
\hypertarget{U+2BDB}{U+2BDB}
&⯛&Pholus&\pholus&pholus&
&
\\
\hypertarget{U+2BDC}{U+2BDC}
&⯜&Nessus&\nessus&nessus&
&
\\
\hypertarget{U+2BDD}{U+2BDD}
&⯝&White Moon Selena&\selena&selena&
&
\\
\hypertarget{U+2BDE}{U+2BDE}
&⯞&Black Diamond On Cross&\blackdiamondoncross&blackdiamondoncross&
&
\\
\hypertarget{U+2BDF}{U+2BDF}
&⯟&True Light Moon Arta&\truelightmoonarta&truelightmoonarta&
&
\\
\hypertarget{U+2BE0}{U+2BE0}
&⯠&Cupido&\cupido&cupido&
&
\\
\hypertarget{U+2BE1}{U+2BE1}
&⯡&Hades&\hades&hades&
&
\\
\hypertarget{U+2BE2}{U+2BE2}
&⯢&Zeus&\zeus&zeus&
&
\\
\hypertarget{U+2BE3}{U+2BE3}
&⯣&Kronos&\kronos&kronos&
&
\\
\hypertarget{U+2BE4}{U+2BE4}
&⯤&Apollon&\apollon&apollon&
&
\\
\hypertarget{U+2BE5}{U+2BE5}
&⯥&Admetos&\admetos&admetos&
&
\\
\hypertarget{U+2BE6}{U+2BE6}
&⯦&Vulcanus&\vulcanus&vulcanus&
&
\\
\end{tabular}
\par
\begin{tabular}{@{}
  >{\ttfamily}l@{}
  >{\sffamily}c@{\hspace{2mm}}
  >{\sffamily\raggedright\vstrut}p{.3\textwidth}
  >{\Large\symbats}c
  >{\ttfamily\raggedright\arraybackslash\textbackslash}m{.29\textwidth}
  >{\ttfamily}l@{\hspace{2mm}}
  >{\ttfamily}l@{}}
\normalfont\bfseries Code&
\normalfont\bfseries Char&
\normalfont\bfseries Unicode name&
\normalfont\normalsize\bfseries Glyph&
\multicolumn1{>{\normalfont\bfseries}l}{\LaTeX\ name}&
\normalfont\bfseries Dupes&
\normalfont\bfseries Compl\\\hline\vrule height1.75em width0pt
\hypertarget{U+2BE7}{U+2BE7}
&⯧&Poseidon&\poseidon&poseidon&
&
\\
\end{tabular}

\clearpage
\section{Symbats3 in alphabetical order of \LaTeX\ name}\label{latexname}
\begin{tabular}{@{}
  >{\ttfamily}l@{}
  >{\sffamily}c@{\hspace{2mm}}
  >{\sffamily\raggedright\vstrut}p{.3\textwidth}
  >{\Large\symbats}c
  >{\ttfamily\raggedright\arraybackslash\textbackslash}m{.29\textwidth}
  >{\ttfamily}l@{\hspace{2mm}}
  >{\ttfamily}l@{}}
\normalfont\bfseries Code&
\normalfont\bfseries Char&
\normalfont\bfseries Unicode name&
\normalfont\normalsize\bfseries Glyph&
\multicolumn1{>{\normalfont\bfseries}l}{\LaTeX\ name}&
\normalfont\bfseries Dupes&
\normalfont\bfseries Compl\\\hline\vrule height1.75em width0pt
U+2BE5
&⯥&Admetos&\admetos&admetos&
&
\\
U+006B
&k&Latin Small Letter K&\aegishjalmur&aegishjalmur&
&
\\
U+005D
&]&Right Square Bracket&\airelement&airelement&
&
\\
U+2220
&∠&Angle&\angleacute&angleacute&
&
\\
U+2625
&☥&Ankh&\ankh&ankh&
\hyperlink{U+00E0}{U+00E0}
&
\\
U+2BE4
&⯤&Apollon&\apollon&apollon&
&
\\
U+2652
&♒&Aquarius&\aquarius&aquarius&
\hyperlink{U+2013}{U+2013}
&
\\
U+2648
&♈&Aries&\aries&aries&
\hyperlink{U+00A1}{U+00A1}
&
\\
U+260A
&☊&Ascending Node&\ascendingnode&ascendingnode&
\hyperlink{U+0058}{U+0058}
&
\\
U+2BD9
&⯙&Astraea&\astraea&astraea&
&
\\
U+0064
&d&Latin Small Letter D&\athameeast&athameeast&
&
\\
U+0073
&s&Latin Small Letter S&\athamenorth&athamenorth&
&
\\
U+0061
&a&Latin Small Letter A&\athamesouth&athamesouth&
&
\\
U+0066
&f&Latin Small Letter F&\athamewest&athamewest&
&
\\
U+0068
&h&Latin Small Letter H&\awen&awen&
&
\\
U+00CC
&Ì&Latin Capital Letter I with Grave&\beltane&beltane&
&
\\
U+2BDE
&⯞&Black Diamond On Cross&\blackdiamondoncross&blackdiamondoncross&
&
\\
U+26B8
&⚸&Black Moon Lilith&\blackmoonlilith&blackmoonlilith&
&
\\
U+264B
&♋&Cancer&\cancer&cancer&
\hyperlink{U+00A2}{U+00A2}
&
\\
U+2651
&♑&Capricorn&\capricorn&capricorn&
\hyperlink{U+00BA}{U+00BA}
&
\\
U+003F
&?&Question Mark&\celtictrefoilknotopen&celtictrefoilknotopen&
&
\\
U+007C
&\|&Vertical Line&\celtictrefoilknotsolid&celtictrefoilknotsolid&
&
\\
U+26B3
&⚳&Ceres&\ceres&ceres&
\hyperlink{U+002C}{U+002C}
&
\\
U+26B7
&⚷&Chiron&\chiron&chiron&
&
\\
U+2295
&⊕&Circled Plus&\circleplus&circleplus&
&
\\
U+260C
&☌&Conjunction&\conjunction&conjunction&
\hyperlink{U+0043}{U+0043}
&
\\
U+002F
&/&Solidus&\crossedstaff&crossedstaff&
&
\\
U+2BE0
&⯠&Cupido&\cupido&cupido&
&
\\
U+260B
&☋&Descending Node&\descendingnode&descendingnode&
\hyperlink{U+005A}{U+005A}
&
\\
U+26A2
&⚢&Doubled Female Sign&\doubledfemale&doubledfemale&
&
\\
U+26A3
&⚣&Doubled Male Sign&\doubledmale&doubledmale&
&
\\
U+2609
&☉&Sun&\dupsun&dupsun&
&
\\
U+004D
&M&Latin Capital Letter M&\dupsunfourspokes&dupsunfourspokes&
&
\\
U+007D
&\}&Right Curly Bracket&\earthelement&earthelement&
&
\\
U+2641
&♁&Earth&\earth&earth&
\hyperlink{U+003C}{U+003C}
&
\\
U+006C
&l&Latin Small Letter L&\eightsabbats&eightsabbats&
&
\\
\end{tabular}
\par
\begin{tabular}{@{}
  >{\ttfamily}l@{}
  >{\sffamily}c@{\hspace{2mm}}
  >{\sffamily\raggedright\vstrut}p{.3\textwidth}
  >{\Large\symbats}c
  >{\ttfamily\raggedright\arraybackslash\textbackslash}m{.29\textwidth}
  >{\ttfamily}l@{\hspace{2mm}}
  >{\ttfamily}l@{}}
\normalfont\bfseries Code&
\normalfont\bfseries Char&
\normalfont\bfseries Unicode name&
\normalfont\normalsize\bfseries Glyph&
\multicolumn1{>{\normalfont\bfseries}l}{\LaTeX\ name}&
\normalfont\bfseries Dupes&
\normalfont\bfseries Compl\\\hline\vrule height1.75em width0pt
U+00CF
&Ï&Latin Capital Letter I with Diaeresis&\eostre&eostre&
&
\\
U+2640
&♀&Female Sign&\female&female&
&
\\
U+005B
&[&Left Square Bracket&\fireelement&fireelement&
&
\\
U+006A
&j&Latin Small Letter J&\fylfotilkley&fylfotilkley&
&
\\
U+264A
&♊&Gemini&\gemini&gemini&
\hyperlink{U+0023}{U+0023}
&
\\
U+00F2
&ò&Latin Small Letter O with Grave&\goddessasmoon&goddessasmoon&
&
\\
U+2BE1
&⯡&Hades&\hades&hades&
&
\\
U+0067
&g&Latin Small Letter G&\hammer&hammer&
&
\\
U+26A9
&⚩&Horizontal Male with Stroke Sign&\horizontalmalewithstroke&horizontalmalewithstroke&
&
\\
U+2BDA
&⯚&Hygiea&\hygiea&hygiea&
&
\\
U+00CE
&Î&Latin Capital Letter I with Circumflex&\imbolc&imbolc&
&
\\
U+26A4
&⚤&Interlocked Female and Male Sign&\interlockedmalefemale&interlockedmalefemale&
&
\\
U+26B5
&⚵&Juno&\juno&juno&
&
\\
U+2643
&♃&Jupiter&\jupiter&jupiter&
\hyperlink{U+0063}{U+0063}
&
\\
U+2BE3
&⯣&Kronos&\kronos&kronos&
&
\\
U+004C
&L&Latin Capital Letter L&\labrys&labrys&
&
\\
U+264C
&♌&Leo&\leo&leo&
\hyperlink{U+221E}{U+221E}
&
\\
U+264E
&♎&Libra&\libra&libra&
\hyperlink{U+00B6}{U+00B6}
&
\\
U+0027
&'&Apostrophe&\lightning&lightning&
&
\\
U+00D4
&Ô&Latin Capital Letter O with Circumflex&\lughnasadh&lughnasadh&
&
\\
U+2642
&♂&Male Sign&\male&male&
&
\\
U+26A5
&⚥&Male and Female Sign&\maleandfemale&maleandfemale&
&
\\
U+26A6
&⚦&Male with Stroke Sign&\malewithstroke&malewithstroke&
&
\\
U+26A7
&⚧&Male with Stroke and Male and Female Sign&\malewithstrokemalefemale&malewithstrokemalefemale&
&
\\
U+0078
&x&Latin Small Letter X&\mars&mars&
&
\\
U+0048
&H&Latin Capital Letter H&\mazeround&mazeround&
&
\\
U+0047
&G&Latin Capital Letter G&\mazesquarebase&mazesquarebase&
&
\\
U+263F
&☿&Mercury&\mercury&mercury&
\hyperlink{U+0060}{U+0060}
&
\\
U+0024
&\$&Dollar Sign&\moonfirstquarteroutline&moonfirstquarteroutline&
&
\\
U+0028
&(&Left Parenthesis&\moonfirstquartersolid&moonfirstquartersolid&
&
\\
U+0029
&)&Right Parenthesis&\moonfull&moonfull&
&
\\
U+005E
&\^{}&Circumflex Accent&\moonlastquarteroutline&moonlastquarteroutline&
&
\\
U+005F
&\_&Low Line&\moonlastquartersolid&moonlastquartersolid&
&
\\
U+0025
&\%&Percent Sign&\moonnew&moonnew&
&
\\
U+263E
&☾&Last Quarter Moon&\moonwaningoutline&moonwaningoutline&
\hyperlink{U+0026}{U+0026}
&
\\
U+002B
&+&Plus Sign&\moonwaningsolid&moonwaningsolid&
&
\\
\end{tabular}
\par
\begin{tabular}{@{}
  >{\ttfamily}l@{}
  >{\sffamily}c@{\hspace{2mm}}
  >{\sffamily\raggedright\vstrut}p{.3\textwidth}
  >{\Large\symbats}c
  >{\ttfamily\raggedright\arraybackslash\textbackslash}m{.29\textwidth}
  >{\ttfamily}l@{\hspace{2mm}}
  >{\ttfamily}l@{}}
\normalfont\bfseries Code&
\normalfont\bfseries Char&
\normalfont\bfseries Unicode name&
\normalfont\normalsize\bfseries Glyph&
\multicolumn1{>{\normalfont\bfseries}l}{\LaTeX\ name}&
\normalfont\bfseries Dupes&
\normalfont\bfseries Compl\\\hline\vrule height1.75em width0pt
U+0021
&!&Exclamation Mark&\moonwaxfullwaneoutline&moonwaxfullwaneoutline&
&
\\
U+0040
&@&Commercial At&\moonwaxfullwanesolid&moonwaxfullwanesolid&
&
\\
U+263D
&☽&First Quarter Moon&\moonwaxingoutline&moonwaxingoutline&
\hyperlink{U+00A3}{U+00A3}
&
\\
U+002A
&*&Asterisk&\moonwaxingsolid&moonwaxingsolid&
&
\\
U+2BC9
&⯉&Neptune Form Two&\neptuneformtwo&neptuneformtwo&
&
\\
U+2646
&♆&Neptune&\neptune&neptune&
\hyperlink{U+006E}{U+006E}
&
\\
U+2BDC
&⯜&Nessus&\nessus&nessus&
&
\\
U+26B2
&⚲&Neuter&\neuter&neuter&
&
\\
U+0034
&4&Digit Four&\octile&octile&
&
\\
U+1690
&ᚐ&Ogham Letter Ailm&\OGailm&OGailm&
&
\\
U+1681
&ᚁ&Ogham Letter Beith&\OGbeith&OGbeith&
&
\\
U+168A
&ᚊ&Ogham Letter Ceirt&\OGceirt&OGceirt&
&
\\
U+1689
&ᚉ&Ogham Letter Coll&\OGcoll&OGcoll&
&
\\
U+1687
&ᚇ&Ogham Letter Dair&\OGdair&OGdair&
&
\\
U+1695
&ᚕ&Ogham Letter Eabhadh&\OGeabhadh&OGeabhadh&
&
\\
U+1693
&ᚓ&Ogham Letter Eadhadh&\OGeadhadh&OGeadhadh&
&
\\
U+1699
&ᚙ&Ogham Letter Eamhancholl&\OGeamhancholl&OGeamhancholl&
&
\\
U+1683
&ᚃ&Ogham Letter Fearn&\OGfearn&OGfearn&
&
\\
U+169B
&᚛&Ogham Feather Mark&\OGfeather&OGfeather&
&
\\
U+168C
&ᚌ&Ogham Letter Gort&\OGgort&OGgort&
&
\\
U+1698
&ᚘ&Ogham Letter Ifin&\OGifin&OGifin&
&
\\
U+1694
&ᚔ&Ogham Letter Iodhadh&\OGiodhadh&OGiodhadh&
&
\\
U+1682
&ᚂ&Ogham Letter Luis&\OGluis&OGluis&
&
\\
U+168B
&ᚋ&Ogham Letter Muin&\OGmuin&OGmuin&
&
\\
U+168D
&ᚍ&Ogham Letter Ngeadal&\OGngeadal&OGngeadal&
&
\\
U+1685
&ᚅ&Ogham Letter Nion&\OGnion&OGnion&
&
\\
U+1691
&ᚑ&Ogham Letter Onn&\OGonn&OGonn&
&
\\
U+1696
&ᚖ&Ogham Letter Or&\OGor&OGor&
&
\\
U+169A
&ᚚ&Ogham Letter Peith&\OGpeith&OGpeith&
&
\\
U+169C
&᚜&Ogham Reversed Feather Mark&\OGreversedfeather&OGreversedfeather&
&
\\
U+168F
&ᚏ&Ogham Letter Ruis&\OGruis&OGruis&
&
\\
U+1684
&ᚄ&Ogham Letter Sail&\OGsail&OGsail&
&
\\
U+168E
&ᚎ&Ogham Letter Straif&\OGstraif&OGstraif&
&
\\
U+1688
&ᚈ&Ogham Letter Tinne&\OGtinne&OGtinne&
&
\\
U+1686
&ᚆ&Ogham Letter Uath&\OGuath&OGuath&
&
\\
U+1697
&ᚗ&Ogham Letter Uilleann&\OGuilleann&OGuilleann&
&
\\
\end{tabular}
\par
\begin{tabular}{@{}
  >{\ttfamily}l@{}
  >{\sffamily}c@{\hspace{2mm}}
  >{\sffamily\raggedright\vstrut}p{.3\textwidth}
  >{\Large\symbats}c
  >{\ttfamily\raggedright\arraybackslash\textbackslash}m{.29\textwidth}
  >{\ttfamily}l@{\hspace{2mm}}
  >{\ttfamily}l@{}}
\normalfont\bfseries Code&
\normalfont\bfseries Char&
\normalfont\bfseries Unicode name&
\normalfont\normalsize\bfseries Glyph&
\multicolumn1{>{\normalfont\bfseries}l}{\LaTeX\ name}&
\normalfont\bfseries Dupes&
\normalfont\bfseries Compl\\\hline\vrule height1.75em width0pt
U+1692
&ᚒ&Ogham Letter Ur&\OGur&OGur&
&
\\
U+00E0
&à&Latin Small Letter A with Grave&\oldankh&oldankh&
\hyperlink{U+2625}{U+2625}
&
\\
U+2013
&–&En Dash&\oldaquarius&oldaquarius&
\hyperlink{U+2652}{U+2652}
&
\\
U+00A1
&¡&Inv Exclamation Mark&\oldaries&oldaries&
\hyperlink{U+2648}{U+2648}
&
\\
U+0058
&X&Latin Capital Letter X&\oldascendingnode&oldascendingnode&
\hyperlink{U+260A}{U+260A}
&
\\
U+00A2
&¢&Cent Sign&\oldcancer&oldcancer&
\hyperlink{U+264B}{U+264B}
&
\\
U+00BA
&º&Masculine Ordinal Indicator&\oldcapricorn&oldcapricorn&
\hyperlink{U+2651}{U+2651}
&
\\
U+002C
&\,&Comma&\oldceres&oldceres&
\hyperlink{U+26B3}{U+26B3}
&
\\
U+0043
&C&Latin Capital Letter C&\oldconjunction&oldconjunction&
\hyperlink{U+260C}{U+260C}
&
\\
U+005A
&Z&Latin Capital Letter Z&\olddescendingnode&olddescendingnode&
\hyperlink{U+260B}{U+260B}
&
\\
U+003C
&\<&Less-Than Sign&\oldearth&oldearth&
\hyperlink{U+2641}{U+2641}
&
\\
U+0023
&\#&Number Sign&\oldgemini&oldgemini&
\hyperlink{U+264A}{U+264A}
&
\\
U+0063
&c&Latin Small Letter C&\oldjupiter&oldjupiter&
\hyperlink{U+2643}{U+2643}
&
\\
U+221E
&∞&Infinity&\oldleo&oldleo&
\hyperlink{U+264C}{U+264C}
&
\\
U+00B6
&¶&Pilcrow Sign&\oldlibra&oldlibra&
\hyperlink{U+264E}{U+264E}
&
\\
U+0060
&\`{}&Grave Accent&\oldmercury&oldmercury&
\hyperlink{U+263F}{U+263F}
&
\\
U+0026
&\&&Ampersand&\oldmoonwaningoutline&oldmoonwaningoutline&
\hyperlink{U+263E}{U+263E}
&
\\
U+00A3
&£&Pound Sign&\oldmoonwaxingoutline&oldmoonwaxingoutline&
\hyperlink{U+263D}{U+263D}
&
\\
U+006E
&n&Latin Small Letter N&\oldneptune&oldneptune&
\hyperlink{U+2646}{U+2646}
&
\\
U+0056
&V&Latin Capital Letter V&\oldopposition&oldopposition&
\hyperlink{U+260D}{U+260D}
&
\\
U+0155
&ŕ&Latin Small Letter R with Acute&\oldpentagramintL&oldpentagramintL&
\hyperlink{U+26E6}{U+26E6}
&
\\
U+0069
&i&Latin Small Letter I&\oldpentagramintRsolid&oldpentagramintRsolid&
\hyperlink{U+00ED}{U+00ED}
&
\\
U+0072
&r&Latin Small Letter R&\oldpentagramintR&oldpentagramintR&
\hyperlink{U+26E5}{U+26E5}
&
\hyperlink{U+0155}{U+0155}
\\
U+0055
&U&Latin Capital Letter U&\oldpentagraminvsolid&oldpentagraminvsolid&
\hyperlink{U+00CD}{U+00CD}
&
\\
U+0057
&W&Latin Capital Letter W&\oldpentagraminv&oldpentagraminv&
\hyperlink{U+26E7}{U+26E7}
&
\\
U+0077
&w&Latin Small Letter W&\oldpentagram&oldpentagram&
\hyperlink{U+26E4}{U+26E4}
&
\\
U+2260
&≠&Not Equal To&\oldpisces&oldpisces&
\hyperlink{U+2653}{U+2653}
&
\\
U+006D
&m&Latin Small Letter M&\oldpluto&oldpluto&
\hyperlink{U+2647}{U+2647}
&
\\
U+00AA
&ª&Feminine Ordinal Indicator&\oldsaggitarius&oldsaggitarius&
\hyperlink{U+2650}{U+2650}
&
\\
U+0076
&v&Latin Small Letter V&\oldsaturn&oldsaturn&
\hyperlink{U+2644}{U+2644}
&
\\
U+0042
&B&Latin Capital Letter B&\oldsextile&oldsextile&
\hyperlink{U+26B9}{U+26B9}
&
\\
U+004E
&N&Latin Capital Letter N&\oldsquare&oldsquare&
\hyperlink{U+25A1}{U+25A1}
&
\\
U+0033
&3&Digit Three&\oldsuncoronafourspokes&oldsuncoronafourspokes&
\hyperlink{U+0038}{U+0038}
&
\\
U+2122
&™&Trade Mark Sign&\oldtaurus&oldtaurus&
\hyperlink{U+2649}{U+2649}
&
\\
U+0062
&b&Latin Small Letter B&\olduranus&olduranus&
\hyperlink{U+2645}{U+2645}
&
\\
U+00A7
&§&Section Sign&\oldvirgo&oldvirgo&
\hyperlink{U+264D}{U+264D}
&
\\
\end{tabular}
\par
\begin{tabular}{@{}
  >{\ttfamily}l@{}
  >{\sffamily}c@{\hspace{2mm}}
  >{\sffamily\raggedright\vstrut}p{.3\textwidth}
  >{\Large\symbats}c
  >{\ttfamily\raggedright\arraybackslash\textbackslash}m{.29\textwidth}
  >{\ttfamily}l@{\hspace{2mm}}
  >{\ttfamily}l@{}}
\normalfont\bfseries Code&
\normalfont\bfseries Char&
\normalfont\bfseries Unicode name&
\normalfont\normalsize\bfseries Glyph&
\multicolumn1{>{\normalfont\bfseries}l}{\LaTeX\ name}&
\normalfont\bfseries Dupes&
\normalfont\bfseries Compl\\\hline\vrule height1.75em width0pt
U+260D
&☍&Opposition&\opposition&opposition&
\hyperlink{U+0056}{U+0056}
&
\\
U+26B4
&⚴&Pallas&\pallas&pallas&
&
\\
U+00C8
&È&Latin Capital Letter E with Grave&\parting&parting&
&
\\
U+0071
&q&Latin Small Letter Q&\pentagramcircled&pentagramcircled&
&
\hyperlink{U+0051}{U+0051}
\\
U+00E9
&é&Latin Small Letter E with Acute&\pentagramintLcircled&pentagramintLcircled&
&
\hyperlink{U+0065}{U+0065}
\\
U+26E6
&⛦&Left-Handed Interlaced Pentagram&\pentagramintL&pentagramintL&
\hyperlink{U+0155}{U+0155}
&
\\
U+26E5
&⛥&Right-Handed Interlaced Pentagram&\pentagramintR&pentagramintR&
&
\\
U+0065
&e&Latin Small Letter E&\pentagramintRcircled&pentagramintRcircled&
&
\hyperlink{U+00E9}{U+00E9}
\\
U+00ED
&í&Latin Small Letter I with Acute&\pentagramintRsolid&pentagramintRsolid&
\hyperlink{U+0069}{U+0069}
&
\\
U+0050
&P&Latin Capital Letter P&\pentagramintRtriangle&pentagramintRtriangle&
&
\\
U+0051
&Q&Latin Capital Letter Q&\pentagraminvcircled&pentagraminvcircled&
&
\hyperlink{U+0071}{U+0071}
\\
U+0045
&E&Latin Capital Letter E&\pentagraminvintLcircled&pentagraminvintLcircled&
&
\hyperlink{U+00C9}{U+00C9}
\\
U+0052
&R&Latin Capital Letter R&\pentagraminvintL&pentagraminvintL&
&
\hyperlink{U+0154}{U+0154}
\\
U+00C9
&É&Latin Capital Letter E with Acute&\pentagraminvintRcircled&pentagraminvintRcircled&
&
\hyperlink{U+0045}{U+0045}
\\
U+0049
&I&Latin Capital Letter I&\pentagraminvintRsolid&pentagraminvintRsolid&
&
\hyperlink{none}{none}
\\
U+0154
&Ŕ&Latin Capital Letter R with Acute&\pentagraminvintR&pentagraminvintR&
&
\hyperlink{U+0052}{U+0052}
\\
U+00CD
&Í&Latin Capital Letter I with Acute&\pentagraminvsolid&pentagraminvsolid&
\hyperlink{U+0055}{U+0055}
&
\\
U+26E7
&⛧&Inv Pentagram&\pentagraminv&pentagraminv&
\hyperlink{U+0057}{U+0057}
&
\\
U+0079
&y&Latin Small Letter Y&\pentagramrough&pentagramrough&
&
\\
U+0074
&t&Latin Small Letter T&\pentagramroughcircled&pentagramroughcircled&
&
\\
U+0059
&Y&Latin Capital Letter Y&\pentagramroughinv&pentagramroughinv&
&
\\
U+0054
&T&Latin Capital Letter T&\pentagramroughinvcircled&pentagramroughinvcircled&
&
\\
U+004F
&O&Latin Capital Letter O&\pentagramroughinvsolid&pentagramroughinvsolid&
&
\hyperlink{U+006F}{U+006F}
\\
U+006F
&o&Latin Small Letter O&\pentagramroughsolid&pentagramroughsolid&
&
\hyperlink{U+004F}{U+004F}
\\
U+0075
&u&Latin Small Letter U&\pentagramsolid&pentagramsolid&
&
\\
U+26E4
&⛤&Pentagram&\pentagram&pentagram&
\hyperlink{U+0077}{U+0077}
&
\\
U+0070
&p&Latin Small Letter P&\pentagramwithtriangle&pentagramwithtriangle&
&
\\
U+00F9
&ù&Latin Small Letter U with Grave&\perfectcouple&perfectcouple&
&
\\
U+00D9
&Ù&Latin Capital Letter U with Grave&\perfectcouplethirddegree&perfectcouplethirddegree&
&
\\
U+2BDB
&⯛&Pholus&\pholus&pholus&
&
\\
U+2653
&♓&Pisces&\pisces&pisces&
\hyperlink{U+2260}{U+2260}
&
\\
U+2BD6
&⯖&Pluto Form Five&\plutoformfive&plutoformfive&
&
\\
U+2BD5
&⯕&Pluto Form Four&\plutoformfour&plutoformfour&
&
\\
U+2BD4
&&Pluto Form Three&\plutoformthree&plutoformthree&
&
\\
U+2BD3
&⯓&Pluto Form Two&\plutoformtwo&plutoformtwo&
&
\\
U+2647
&♇&Pluto&\pluto&pluto&
\hyperlink{U+006D}{U+006D}
&
\\
\end{tabular}
\par
\begin{tabular}{@{}
  >{\ttfamily}l@{}
  >{\sffamily}c@{\hspace{2mm}}
  >{\sffamily\raggedright\vstrut}p{.3\textwidth}
  >{\Large\symbats}c
  >{\ttfamily\raggedright\arraybackslash\textbackslash}m{.29\textwidth}
  >{\ttfamily}l@{\hspace{2mm}}
  >{\ttfamily}l@{}}
\normalfont\bfseries Code&
\normalfont\bfseries Char&
\normalfont\bfseries Unicode name&
\normalfont\normalsize\bfseries Glyph&
\multicolumn1{>{\normalfont\bfseries}l}{\LaTeX\ name}&
\normalfont\bfseries Dupes&
\normalfont\bfseries Compl\\\hline\vrule height1.75em width0pt
U+2BE7
&⯧&Poseidon&\poseidon&poseidon&
&
\\
U+00C0
&À&Latin Capital Letter A with Grave&\powergoingforth&powergoingforth&
&
\\
U+211E
&℞&Prescription Take (Recipe)&\prescription&prescription&
&
\\
U+2BD8
&⯘&Proserpina&\proserpina&proserpina&
&
\\
U+26BB
&⚻&Quincunx&\quincunx&quincunx&
&
\\
U+003E
&\>&Greater-Than Sign&\recipe&recipe&
&
\\
U+004A
&J&Latin Capital Letter J&\rockartgoddess&rockartgoddess&
&
\\
U+004B
&K&Latin Capital Letter K&\rockarthorse&rockarthorse&
&
\\
U+16AA
&ᚪ&Runic Letter Ac A&\RUaca&RUaca&
&
\\
U+16AB
&ᚫ&Runic Letter Aesc&\RUaesc&RUaesc&
&
\\
U+16C9
&ᛉ&Runic Letter Algiz Eolhx&\RUalgizeolhx&RUalgizeolhx&
&
\\
U+16A8
&ᚨ&Runic Letter Ansuz A&\RUansuza&RUansuza&
&
\\
U+16EE
&ᛮ&Runic Arlaug Symbol&\RUarlaugsymbol&RUarlaugsymbol&
&
\\
U+16F0
&ᛰ&Runic Belgthor Symbol&\RUbelgthorsymbol&RUbelgthorsymbol&
&
\\
U+16D2
&ᛒ&Runic Letter Berkanan Beorc Bjarkan B&\RUberkananbeorcbjarkanb&RUberkananbeorcbjarkanb&
&
\\
U+16CD
&ᛍ&Runic Letter C&\RUc&RUc&
&
\\
U+16E3
&ᛣ&Runic Letter Calc&\RUcalc&RUcalc&
&
\\
U+16E4
&ᛤ&Runic Letter Cealc&\RUcealc&RUcealc&
&
\\
U+16B3
&ᚳ&Runic Letter Cen&\RUcen&RUcen&
&
\\
U+16ED
&᛭&Runic Cross Punctuation&\RUcrosspunctuation&RUcrosspunctuation&
&
\\
U+16E2
&ᛢ&Runic Letter Cweorth&\RUcweorth&RUcweorth&
&
\\
U+16D1
&ᛑ&Runic Letter D&\RUd&RUd&
&
\\
U+16DE
&ᛞ&Runic Letter Dagaz Daeg D&\RUdagazdaegd&RUdagazdaegd&
&
\\
U+16DB
&ᛛ&Runic Letter Dotted-L&\RUdottedl&RUdottedl&
&
\\
U+16C0
&ᛀ&Runic Letter Dotted-N&\RUdottedn&RUdottedn&
&
\\
U+16D4
&ᛔ&Runic Letter Dotted-P&\RUdottedp&RUdottedp&
&
\\
U+16C2
&ᛂ&Runic Letter E&\RUe&RUe&
&
\\
U+16E0
&ᛠ&Runic Letter Ear&\RUear&RUear&
&
\\
U+16D6
&ᛖ&Runic Letter Ehwaz Eh E&\RUehwazehe&RUehwazehe&
&
\\
U+16B6
&ᚶ&Runic Letter Eng&\RUeng&RUeng&
&
\\
U+16A7
&ᚧ&Runic Letter Eth&\RUeth&RUeth&
&
\\
U+16A0
&ᚠ&Runic Letter Fehu Feoh Fe F&\RUfehufeohfef&RUfehufeohfef&
&
\\
U+16F7
&ᛷ&Runic Letter Franks Casket Ac&\RUfrankscasketac&RUfrankscasketac&
&
\\
U+16F8
&ᛸ&Runic Letter Franks Casket Aesc&\RUfrankscasketaesc&RUfrankscasketaesc&
&
\\
U+16F6
&ᛶ&Runic Letter Franks Casket Eh&\RUfrankscasketeh&RUfrankscasketeh&
&
\\
U+16F5
&ᛵ&Runic Letter Franks Casket Is&\RUfrankscasketis&RUfrankscasketis&
&
\\
\end{tabular}
\par
\begin{tabular}{@{}
  >{\ttfamily}l@{}
  >{\sffamily}c@{\hspace{2mm}}
  >{\sffamily\raggedright\vstrut}p{.3\textwidth}
  >{\Large\symbats}c
  >{\ttfamily\raggedright\arraybackslash\textbackslash}m{.29\textwidth}
  >{\ttfamily}l@{\hspace{2mm}}
  >{\ttfamily}l@{}}
\normalfont\bfseries Code&
\normalfont\bfseries Char&
\normalfont\bfseries Unicode name&
\normalfont\normalsize\bfseries Glyph&
\multicolumn1{>{\normalfont\bfseries}l}{\LaTeX\ name}&
\normalfont\bfseries Dupes&
\normalfont\bfseries Compl\\\hline\vrule height1.75em width0pt
U+16F4
&ᛴ&Runic Letter Franks Casket Os&\RUfrankscasketos&RUfrankscasketos&
&
\\
U+16B5
&ᚵ&Runic Letter G&\RUg&RUg&
&
\\
U+16B8
&ᚸ&Runic Letter Gar&\RUgar&RUgar&
&
\\
U+16B7
&ᚷ&Runic Letter Gebo Gyfu G&\RUgebogyfug&RUgebogyfug&
&
\\
U+16C4
&ᛄ&Runic Letter Ger&\RUger&RUger&
&
\\
U+16BB
&ᚻ&Runic Letter Haegl H&\RUhaeglh&RUhaeglh&
&
\\
U+16BA
&ᚺ&Runic Letter Haglaz H&\RUhaglazh&RUhaglazh&
&
\\
U+16E8
&ᛨ&Runic Letter Icelandic-Yr&\RUicelandicyr&RUicelandicyr&
&
\\
U+16DD
&ᛝ&Runic Letter Ing&\RUing&RUing&
&
\\
U+16DC
&ᛜ&Runic Letter Ingwaz&\RUingwaz&RUingwaz&
&
\\
U+16E1
&ᛡ&Runic Letter Ior&\RUior&RUior&
&
\\
U+16C1
&ᛁ&Runic Letter Isaz Is Iss I&\RUisazisissi&RUisazisissi&
&
\\
U+16C7
&ᛇ&Runic Letter Iwaz Eoh&\RUiwazeoh&RUiwazeoh&
&
\\
U+16C3
&ᛃ&Runic Letter Jeran J&\RUjeranj&RUjeranj&
&
\\
U+16F1
&ᛱ&Runic Letter K&\RUk&RUk&
&
\\
U+16B2
&ᚲ&Runic Letter Kauna&\RUkauna&RUkauna&
&
\\
U+16B4
&ᚴ&Runic Letter Kaun K&\RUkaunk&RUkaunk&
&
\\
U+16DA
&ᛚ&Runic Letter Laukaz Lagu Logr L&\RUlaukazlagulogrl&RUlaukazlagulogrl&
&
\\
U+16C5
&ᛅ&Runic Letter Long-Branch-Ar Ae&\RUlongbrancharae&RUlongbrancharae&
&
\\
U+16BC
&ᚼ&Runic Letter Long-Branch-Hagall H&\RUlongbranchhagallh&RUlongbranchhagallh&
&
\\
U+16D8
&ᛘ&Runic Letter Long-Branch-Madr M&\RUlongbranchmadrm&RUlongbranchmadrm&
&
\\
U+16AC
&ᚬ&Runic Letter Long-Branch-Oss O&\RUlongbranchosso&RUlongbranchosso&
&
\\
U+16E6
&ᛦ&Runic Letter Long-Branch-Yr&\RUlongbranchyr&RUlongbranchyr&
&
\\
U+16D7
&ᛗ&Runic Letter Mannaz Man M&\RUmannazmanm&RUmannazmanm&
&
\\
U+16EC
&᛬&Runic Multiple Punctuation&\RUmultiplepunctuation&RUmultiplepunctuation&
&
\\
U+16BE
&ᚾ&Runic Letter Naudiz Nyd Naud N&\RUnaudiznydnaudn&RUnaudiznydnaudn&
&
\\
U+16AE
&ᚮ&Runic Letter O&\RUo&RUo&
&
\\
U+16AF
&ᚯ&Runic Letter Oe&\RUoe&RUoe&
&
\\
U+16B0
&ᚰ&Runic Letter On&\RUon&RUon&
&
\\
U+16F3
&ᛳ&Runic Letter Oo&\RUoo&RUoo&
&
\\
U+16D5
&ᛕ&Runic Letter Open-P&\RUopenp&RUopenp&
&
\\
U+16A9
&ᚩ&Runic Letter Os O&\RUoso&RUoso&
&
\\
U+16DF
&ᛟ&Runic Letter Othalan Ethel O&\RUothalanethelo&RUothalanethelo&
&
\\
U+16C8
&ᛈ&Runic Letter Pertho Peorth P&\RUperthopeorthp&RUperthopeorthp&
&
\\
U+16E9
&ᛩ&Runic Letter Q&\RUq&RUq&
&
\\
U+16B1
&ᚱ&Runic Letter Raido Rad Reid R&\RUraidoradreidr&RUraidoradreidr&
&
\\
\end{tabular}
\par
\begin{tabular}{@{}
  >{\ttfamily}l@{}
  >{\sffamily}c@{\hspace{2mm}}
  >{\sffamily\raggedright\vstrut}p{.3\textwidth}
  >{\Large\symbats}c
  >{\ttfamily\raggedright\arraybackslash\textbackslash}m{.29\textwidth}
  >{\ttfamily}l@{\hspace{2mm}}
  >{\ttfamily}l@{}}
\normalfont\bfseries Code&
\normalfont\bfseries Char&
\normalfont\bfseries Unicode name&
\normalfont\normalsize\bfseries Glyph&
\multicolumn1{>{\normalfont\bfseries}l}{\LaTeX\ name}&
\normalfont\bfseries Dupes&
\normalfont\bfseries Compl\\\hline\vrule height1.75em width0pt
U+16F2
&ᛲ&Runic Letter Sh&\RUsh&RUsh&
&
\\
U+16C6
&ᛆ&Runic Letter Short-Twig-Ar A&\RUshorttwigara&RUshorttwigara&
&
\\
U+16D3
&ᛓ&Runic Letter Short-Twig-Bjarkan B&\RUshorttwigbjarkanb&RUshorttwigbjarkanb&
&
\\
U+16BD
&ᚽ&Runic Letter Short-Twig-Hagall H&\RUshorttwighagallh&RUshorttwighagallh&
&
\\
U+16D9
&ᛙ&Runic Letter Short-Twig-Madr M&\RUshorttwigmadrm&RUshorttwigmadrm&
&
\\
U+16BF
&ᚿ&Runic Letter Short-Twig-Naud N&\RUshorttwignaudn&RUshorttwignaudn&
&
\\
U+16AD
&ᚭ&Runic Letter Short-Twig-Oss O&\RUshorttwigosso&RUshorttwigosso&
&
\\
U+16CC
&ᛌ&Runic Letter Short-Twig-Sol S&\RUshorttwigsols&RUshorttwigsols&
&
\\
U+16D0
&ᛐ&Runic Letter Short-Twig-Tyr T&\RUshorttwigtyrt&RUshorttwigtyrt&
&
\\
U+16E7
&ᛧ&Runic Letter Short-Twig-Yr&\RUshorttwigyr&RUshorttwigyr&
&
\\
U+16CB
&ᛋ&Runic Letter Sigel Long-Branch-Sol S&\RUsigellongbranchsols&RUsigellongbranchsols&
&
\\
U+16EB
&᛫&Runic Single Punctuation&\RUsinglepunctuation&RUsinglepunctuation&
&
\\
U+16CA
&ᛊ&Runic Letter Sowilo S&\RUsowilos&RUsowilos&
&
\\
U+16E5
&ᛥ&Runic Letter Stan&\RUstan&RUstan&
&
\\
U+16A6
&ᚦ&Runic Letter Thurisaz Thurs Thorn&\RUthurisazthursthorn&RUthurisazthursthorn&
&
\\
U+16CF
&ᛏ&Runic Letter Tiwaz Tir Tyr T&\RUtiwaztirtyrt&RUtiwaztirtyrt&
&
\\
U+16EF
&ᛯ&Runic Tvimadur Symbol&\RUtvimadursymbol&RUtvimadursymbol&
&
\\
U+16A2
&ᚢ&Runic Letter Uruz Ur U&\RUuruzuru&RUuruzuru&
&
\\
U+16A1
&ᚡ&Runic Letter V&\RUv&RUv&
&
\\
U+16A5
&ᚥ&Runic Letter W&\RUw&RUw&
&
\\
U+16B9
&ᚹ&Runic Letter Wunjo Wynn W&\RUwunjowynnw&RUwunjowynnw&
&
\\
U+16EA
&ᛪ&Runic Letter X&\RUx&RUx&
&
\\
U+16A4
&ᚤ&Runic Letter Y&\RUy&RUy&
&
\\
U+16A3
&ᚣ&Runic Letter Yr&\RUyr&RUyr&
&
\\
U+16CE
&ᛎ&Runic Letter Z&\RUz&RUz&
&
\\
U+2650
&♐&Sagittarius&\saggitarius&saggitarius&
\hyperlink{U+00AA}{U+00AA}
&
\\
U+00E8
&è&Latin Small Letter E with Grave&\salute&salute&
&
\\
U+00C5
&Å&Latin Capital Letter A with Ring Above&\samhuinn&samhuinn&
&
\\
U+2644
&♄&Saturn&\saturn&saturn&
\hyperlink{U+0076}{U+0076}
&
\\
U+2022
&•&Bullet&\scorpio&scorpio&
&
\\
U+264F
&♏&Scorpius&\scorpius&scorpius&
&
\\
U+00EC
&ì&Latin Small Letter I with Grave&\scourge&scourge&
&
\\
U+2BDD
&⯝&White Moon Selena&\selena&selena&
&
\\
U+26BA
&⚺&Semisextile&\semisextile&semisextile&
&
\\
U+005C
&\textbackslash&Reverse Solidus&\serpentimpaled&serpentimpaled&
&
\\
U+26BC
&⚼&Sesquiquadrate&\sesquiquadrate&sesquiquadrate&
&
\\
\end{tabular}
\par
\begin{tabular}{@{}
  >{\ttfamily}l@{}
  >{\sffamily}c@{\hspace{2mm}}
  >{\sffamily\raggedright\vstrut}p{.3\textwidth}
  >{\Large\symbats}c
  >{\ttfamily\raggedright\arraybackslash\textbackslash}m{.29\textwidth}
  >{\ttfamily}l@{\hspace{2mm}}
  >{\ttfamily}l@{}}
\normalfont\bfseries Code&
\normalfont\bfseries Char&
\normalfont\bfseries Unicode name&
\normalfont\normalsize\bfseries Glyph&
\multicolumn1{>{\normalfont\bfseries}l}{\LaTeX\ name}&
\normalfont\bfseries Dupes&
\normalfont\bfseries Compl\\\hline\vrule height1.75em width0pt
U+26B9
&⚹&Sextile&\sextile&sextile&
\hyperlink{U+0042}{U+0042}
&
\\
U+25A1
&□&White Square&\square&square&
\hyperlink{U+004F}{U+004F}
&
\\
U+003B
&;&Semicolon&\stareightarrows&stareightarrows&
&
\\
U+002E
&.&Full Stop&\sulphur&sulphur&
&
\\
U+0035
&5&Digit Five&\sun&sun&
&
\\
U+0038
&8&Digit Eight&\suncoronafourspokes&suncoronafourspokes&
\hyperlink{U+0033}{U+0033}
&
\\
U+0039
&9&Digit Nine&\suncoronahubfourspokes&suncoronahubfourspokes&
&
\\
U+003D
&=&Equals Sign&\suncoronasixspokes&suncoronasixspokes&
&
\\
U+0037
&7&Digit Seven&\suneightspokes&suneightspokes&
&
\\
U+0036
&6&Digit Six&\sunfourdoublespokes&sunfourdoublespokes&
&
\\
U+0030
&0&Digit Zero&\sunfourspokes&sunfourspokes&
&
\\
U+263C
&☼&White Sun with Rays&\sunrays&sunrays&
&
\\
U+002D
&-&Hyphen-Minus&\sunthreedoublespokes&sunthreedoublespokes&
&
\\
U+0031
&1&Digit One&\sunwithrays&sunwithrays&
&
\\
U+0032
&2&Digit Two&\sunwithraysfourspokes&sunwithraysfourspokes&
&
\\
U+2649
&♉&Taurus&\taurus&taurus&
\hyperlink{U+2122}{U+2122}
&
\\
U+2BD7
&⯗&Transpluto&\transpluto&transpluto&
&
\\
U+2BDF
&⯟&True Light Moon Arta&\truelightmoonarta&truelightmoonarta&
&
\\
U+2645
&♅&Uranus&\uranus&uranus&
\hyperlink{U+0062}{U+0062}
&
\\
U+003A
&:&Colon&\valknutoutline&valknutoutline&
&
\\
U+0022
&“&Quotation Mark&\valknutsolid&valknutsolid&
&
\\
U+007E
&\~{}&Tilde&\vargemini&vargemini&
&
\\
U+26E2
&⛢&Astronomical Symbol For Uranus&\varuranus&varuranus&
&
\\
U+007A
&z&Latin Small Letter Z&\venus&venus&
&
\\
U+00D3
&Ó&Latin Capital Letter O with Acute&\venusfigureoutline&venusfigureoutline&
&
\\
U+00F3
&ó&Latin Small Letter O with Acute&\venusfiguresolid&venusfiguresolid&
&
\\
U+26A8
&⚨&Vertical Male with Stroke Sign&\verticalmalewithstroke&verticalmalewithstroke&
&
\\
U+26B6
&⚶&Vesta&\vesta&vesta&
&
\\
U+264D
&♍&Virgo&\virgo&virgo&
\hyperlink{U+00A7}{U+00A7}
&
\\
U+2BE6
&⯦&Vulcanus&\vulcanus&vulcanus&
&
\\
U+007B
&\{&Left Curly Bracket&\waterelement&waterelement&
&
\\
U+0044
&D&Latin Capital Letter D&\whiteknifeeast&whiteknifeeast&
&
\\
U+0053
&S&Latin Capital Letter S&\whiteknifenorth&whiteknifenorth&
&
\\
U+0041
&A&Latin Capital Letter A&\whiteknifesouth&whiteknifesouth&
&
\\
U+0046
&F&Latin Capital Letter F&\whiteknifewest&whiteknifewest&
&
\\
U+0132
&IJ&Latin Capital Ligature Ij&\yule&yule&
&
\\
\end{tabular}
\par
\begin{tabular}{@{}
  >{\ttfamily}l@{}
  >{\sffamily}c@{\hspace{2mm}}
  >{\sffamily\raggedright\vstrut}p{.3\textwidth}
  >{\Large\symbats}c
  >{\ttfamily\raggedright\arraybackslash\textbackslash}m{.29\textwidth}
  >{\ttfamily}l@{\hspace{2mm}}
  >{\ttfamily}l@{}}
\normalfont\bfseries Code&
\normalfont\bfseries Char&
\normalfont\bfseries Unicode name&
\normalfont\normalsize\bfseries Glyph&
\multicolumn1{>{\normalfont\bfseries}l}{\LaTeX\ name}&
\normalfont\bfseries Dupes&
\normalfont\bfseries Compl\\\hline\vrule height1.75em width0pt
U+2BE2
&⯢&Zeus&\zeus&zeus&
&
\\
\end{tabular}

\clearpage
\section{Symbats3 in alphabetical order of Unicode name}\label{unicodename}
\input{unicodename}
%\clearpage
%\input{test}
\end{document}
